\documentclass[11pt, a4paper]{article}
\usepackage[utf8]{inputenc}
\usepackage[spanish]{babel}
\usepackage{amsmath, amssymb, amsfonts}
\usepackage{geometry}
\geometry{left=2.5cm, right=2.5cm, top=2.5cm, bottom=2.5cm}
\usepackage{tikz}
\usepackage{bm}

% Línea discontinua entre ejercicios
\newcommand{\separador}{\par\vspace{0.3cm}\noindent\begin{tikzpicture}\draw[dashed] (0,0) -- (\linewidth,0);\end{tikzpicture}\par\vspace{0.3cm}}

\title{\vspace{-1.5cm}\textbf{Resolución de Ecuaciones Diferenciales Ordinarias}\\ \large Métodos Matemáticos I}
\author{\textbf{César Endris Acosta}}

\begin{document}

\maketitle
\vspace{0.2cm}

\section{Hoja 1}

\noindent\textbf{1. Integrad la EDO: $\bm{(1+y^2)dx + (1+x^2)dy = 0}$} \\
\begin{align*}
\frac{dx}{1+x^2} &= -\frac{dy}{1+y^2}
\end{align*}
$y = \tan(\theta) \implies dy = \sec^2(\theta)\,d\theta,\ \ 1+y^2=\sec^2(\theta)$.
\[
\int \frac{dy}{1+y^2} = \int \frac{\sec^2(\theta)\,d\theta}{\sec^2(\theta)} = \int d\theta = \theta = \arctan(y)
\]
\begin{align*}
\arctan(x) &= -\arctan(y) + C
\end{align*}
\textit{Los denominadores $1+x^2$ y $1+y^2$ son siempre positivos ($\geq 1$) $\forall x,y\in\mathbb{R}$, por lo que no existe singularidad ni restricción de signo sobre $C$.}
\[
\boxed{\arctan(x) + \arctan(y) = C}
\]


\noindent\textbf{2. Integrad la EDO: $\bm{(1+y^2)dx + xydy = 0}$} \\
\begin{align*}
\frac{dx}{x} &= -\frac{y}{1+y^2}dy
\end{align*}
$u = 1+y^2 \implies du = 2y\,dy $.
\[
\int \frac{y}{1+y^2}dy = \int \frac{du/2}{u} = \frac{1}{2}\ln|u| = \frac{1}{2}\ln(1+y^2)
\]
\begin{align*}
\int \frac{dx}{x} &= -\frac{1}{2}\int \frac{du}{u} \\
\ln|x| &= -\frac{1}{2}\ln(1+y^2) + C
\end{align*}
\textit{Al dividir por $x$ se excluye $x=0$; sustituyendo en la ecuación original, $(1+y^2)\cdot 0 + 0\cdot y\,dy=0$ se cumple, así que la recta $x=0$ es también solución. En el caso ($x\neq0$), al despejar $x= \pm Ke^{-\frac12\ln(1+y^2)}$ con $K>0$, pero como $x^2(1+y^2)\geq0$ siempre, la constante final debe cumplir $K\geq0$ (con $K=0$ recuperando la solución singular $x=0$).}
\[
\boxed{x^2(1+y^2) = K, \qquad K\geq 0}
\]


\noindent\textbf{3. Integrad la EDO: $\bm{y'\sin(x)-y\cos(x)=0}$ con la condición inicial $\bm{y(\pi/2)=1}$} \\
\begin{align*}
\frac{dy}{y} &= \frac{\cos(x)}{\sin(x)}dx
\end{align*}
\textit{Al dividir por $y$ se excluye $y=0$; pero si $y\equiv 0$, entonces $y'=0$ y la ecuación original $0\cdot\sin(x)-0\cdot\cos(x)=0$ se satisface $\forall x$, luego $y=0$ es también solución. Con $y\neq0$:}
\begin{align*}
\int \frac{dy}{y} &= \int \cot(x)dx \\
\ln|y| &= \ln|\sin(x)| + C
\end{align*}
Al exponenciar, $|y|=e^C|\sin(x)|$, y separando $y>0$ o $y<0$ se obtiene $y=\pm K\sin(x)$ con $K>0$; incluyendo $K=0$ se recupera el caso singular.
Aplicando $y(\pi/2)=1 \implies 1 = K\sin(\pi/2) \implies K = 1$.
\[
\boxed{y = \sin(x)}
\]


\noindent\textbf{4. Integrad la EDO: $\bm{x\sqrt{1+y^2}+yy'\sqrt{1+x^2}=0}$} \\
\begin{align*}
\frac{x}{\sqrt{1+x^2}}dx &= -\frac{y}{\sqrt{1+y^2}}dy
\end{align*}
\[
\int \frac{y}{\sqrt{1+y^2}}dy = \int \frac{du/2}{\sqrt{u}} = \sqrt{u} = \sqrt{1+y^2}
\]
\[
\int \frac{x}{\sqrt{1+x^2}}dx=\sqrt{1+x^2}
\]
\begin{align*}
\sqrt{1+x^2} &= -\sqrt{1+y^2} + C
\end{align*}
\textit{$\sqrt{1+x^2}\geq1$ y $\sqrt{1+y^2}\geq1$ siempre, no hay singularidad, pero sí una restricción: $C\geq2$.}
\[
\boxed{\sqrt{1+x^2} + \sqrt{1+y^2} = C, \qquad C\geq 2}
\]


\noindent\textbf{5. Integrad la EDO: $\bm{x\sqrt{1-y^2}dx+y\sqrt{1-x^2}dy=0}$ con la condición inicial $\bm{y(0)=1}$} \\
\begin{align*}
\frac{x}{\sqrt{1-x^2}}dx &= -\frac{y}{\sqrt{1-y^2}}dy \\
\int \frac{x}{\sqrt{1-x^2}}dx &= -\int \frac{y}{\sqrt{1-y^2}}dy \\
-\sqrt{1-x^2} &= \sqrt{1-y^2} + C \implies \sqrt{1-x^2} + \sqrt{1-y^2} = K
\end{align*}
\textit{El dominio exige $|x|\leq1$, $|y|\leq1$; en $x=\pm1$ o $y=\pm1$ las raíces se anulan (no constituyen curvas de solución independientes). Como $0\leq\sqrt{1-x^2}\leq1$ y $0\leq\sqrt{1-y^2}\leq1$, se requiere $0\leq K\leq2$.}
Aplicando $y(0)=1 \implies \sqrt{1-0} + \sqrt{1-1} = K \implies K=1$.
\[
\boxed{\sqrt{1-x^2} + \sqrt{1-y^2} = 1}
\]


\noindent\textbf{6. Integrad la EDO: $\bm{e^{-y}(1+y')=1}$} \\
\begin{align*}
1+y' &= e^y \implies y' = e^y - 1 \implies \frac{dy}{e^y-1} = dx
\end{align*}
\textit{Al dividir por $e^y-1$ se excluye $y=0$; si $y\equiv0$, entonces $y'=0$ y la ecuación original $e^{0}(1+0)=1$ se cumple, luego $y=0$ es también solución (caso singular). En el caso $y\neq0$:}

$u=e^y-1 \implies du=e^y\,dy=(u+1)\,dy \implies dy=\dfrac{du}{u+1}$.
\begin{align*}
\frac{dy}{e^y-1} &= \frac{du/(u+1)}{u} = \frac{du}{u(u+1)} \implies \dfrac{1}{u(u+1)}=( \dfrac{1}{u}-\dfrac{1}{u+1} ){du}.
\end{align*}

\begin{align*}
\int \left(\frac{1}{u}-\frac{1}{u+1}\right)du &= \int dx \\
\ln\left|\frac{u}{u+1}\right| &= x+C
\end{align*}
Deshaciendo el cambio: $\ln\left|\dfrac{e^y-1}{e^y}\right|=x+C \implies \ln|1-e^{-y}|=x+C$.
\textit{Al exponenciar, $K>0$ (sin embargo con $\pm$ se recuperan todas las soluciones, positivo si $y>0$, negativo si $y<0$) y $K=0$ recupera el caso singular $y=0$.}
\[
\boxed{1 - e^{-y} = \pm K e^x}
\]


\noindent\textbf{7. Integrad la EDO: $\bm{y\ln y \, dx + x \, dy = 0}$ con la condición inicial $\bm{y(1)=1}$} \\
\begin{align*}
\frac{dx}{x} &= -\frac{dy}{y\ln y}
\end{align*}
\textit{El dominio requiere $y>0$; al dividir por $\ln y$ se excluye $y=1$. Si $y\equiv1$, entonces $y'=0$ y la ecuación original $1\cdot\ln(1)\,dx+x\cdot0=0$ se cumple, luego $y=1$ es también solución (caso singular).}

$u = \ln y \implies du = \dfrac{dy}{y}$.
\begin{align*}
\int \frac{dx}{x} &= -\int \frac{du}{u} \\
\ln|x| &= -\ln|u| + C \implies x\ln y = K
\end{align*}
Aplicando $y(1)=1 \implies 1 \cdot \ln(1) = K \implies K = 0$, valor que corresponde al caso hallado arriba.
\[
x\ln y = 0 \implies
\]
\[
\boxed{y = 1}
\]


\noindent\textbf{8. Integrad la EDO: $\bm{y' = a^{x+y}}$ con $\bm{a>0}$ y $\bm{a \neq 1}$} \\
$z = x+y \implies z' = 1+y' \implies y' = z'-1$.
\begin{align*}
z'-1 &= a^z \implies z' = 1+a^z \implies \frac{dz}{1+a^z} = dx
\end{align*}
$u = 1+a^z \implies du = a^z\ln(a)\,dz = (u-1)\ln(a)\,dz \implies dz = \dfrac{du}{(u-1)\ln a}$.
\[
\frac{dz}{1+a^z} = \frac{dz}{u} = \frac{du}{u(u-1)\ln a} = (\dfrac{1}{u-1}-\dfrac{1}{u})(\frac{du}{\ln a})
\]

\begin{align*}
\int \frac{dz}{1+a^z} &= \frac{1}{\ln a}\int \left(\frac{1}{u-1}-\frac{1}{u}\right)du = \frac{1}{\ln a}\ln\left|\frac{u-1}{u}\right| \\
\frac{1}{\ln a}\ln\left|\frac{u-1}{u}\right| &= x+C
\end{align*}
\textit{Como $a^z>0\ \forall z\in\mathbb{R}$, se tiene $u=1+a^z>1$ siempre, luego $u\neq0$ y $u-1=a^z\neq0$: ningún paso excluye soluciones.}

\[
\ln\left|\frac{a^z}{1+a^z}\right| = x\ln a + C_1 \implies \frac{a^z}{1+a^z} = Ka^x
\]
\begin{align*}
a^z &= Ka^x(1+a^z) \implies a^z(1-Ka^x) = Ka^x \implies a^z = \frac{Ka^x}{1-Ka^x}
\end{align*}

\[
a^xa^y = \frac{Ka^x}{1-Ka^x} \implies a^y = \frac{K}{1-Ka^x} \implies a^{-y} = \frac{1}{K}-a^x
\]
\textit{Como $a^x>0$ y $a^{-y}>0$ para todo $x,y\in\mathbb{R}$ (pues $a>0$), la constante $\dfrac1K$ debe ser positiva: no hay caso de singularidad ni de signo negativo.}
\[
\boxed{a^x + a^{-y} = K, \qquad K>0}
\]


\noindent\textbf{9. Integrad la EDO: $\bm{y' = \sin(x-y)}$} \\
$z = x-y \implies z' = 1-y'$.
\[
1-z' = \sin(z) \implies z' = 1-\sin(z) \implies \frac{dz}{1-\sin(z)} = dx
\]

$u=\tan\left(\dfrac{z}{2}\right) \implies du = \dfrac{1+u^2}{2}\,dz \implies dz=\dfrac{2\,du}{1+u^2}$.\quad
\begin{tikzpicture}[baseline={([yshift=-0.6ex]current bounding box.center)}, scale=0.5]
\coordinate (A) at (0,0);
\coordinate (B) at (3,0);
\coordinate (C) at (3,2);
\draw[thick] (A) -- (B) -- (C) -- cycle;
\draw (A) ++(0.35,0) arc (0:33.7:0.35);
\node at (0.62,0.17) {$z$};
\node[below] at (1.5,0) {$1-u^2$};
\node[right] at (3,1) {$2u$};
\node[above, sloped] at (1.15,1) {$1+u^2$};
\draw (2.85,0) -- (2.85,0.15) -- (3,0.15);
\end{tikzpicture}

\begin{align*}
\sin(z) = \frac{2u}{1+u^2} \implies 1-\sin(z) &= \frac{(1-u)^2}{1+u^2} \implies \frac{dz}{1-\sin(z)} = \frac{2\,du}{(1-u)^2} \\
\int \frac{2\,du}{(1-u)^2} &= \int dx \implies \frac{2}{1-u} = x + C
\end{align*}

\textit{Al dividir por $1-\sin(z)$ se excluye $\sin(z)=1$, es decir $z=x-y=\dfrac{\pi}{2}+2k\pi$; sustituyendo, $z'=1-\sin(\pi/2)=0$ concuerda con $z$ constante, luego $y=x-\dfrac{\pi}{2}-2k\pi$ es también solución.}
\[
\boxed{\frac{2}{1-\tan\left(\dfrac{x-y}{2}\right)} = x + C}
\]


\noindent\textbf{10. Integrad la EDO: $\bm{y' = ax+by+c}$ con $\bm{a, b, c}$ constantes} \\
Cambio de variable: $z = ax+by+c \implies z' = a+by' \implies y' = \dfrac{z'-a}{b}$.
\begin{align*}
\frac{z'-a}{b} &= z \implies z' = a+bz \implies \frac{dz}{a+bz} = dx
\end{align*}
\textit{Con $b\neq0$ al dividir por $a+bz$ se excluye $z=-a/b$; en ese punto $z'=a+b(-a/b)=0$, consistente con $z$ constante, luego es también solución.}
\begin{align*}
\int \frac{dz}{a+bz} &= \int dx \\
\frac{1}{b}\ln|a+bz| &= x + C_1 \implies a+bz = \pm K e^{bx}
\end{align*}
Se puede tomar cualquier signo según el signo de $a+bz$ será el elegido en el $\pm$ ; $K=0$ recupera el caso singular.
\[
\boxed{a+b(ax+by+c) = \pm K e^{bx}}
\]


\noindent\textbf{11. Integrad la EDO: $\bm{y+xy' = a(1+xy)}$ con la condición inicial $\bm{y(1/a)=-a}$} \\
$z = xy \implies z' = y+xy'$.
\begin{align*}
z' &= a(1+z) \implies \frac{dz}{1+z} = a\,dx
\end{align*}
\textit{Al dividir por $1+z$ se excluye $z=-1$, es decir $xy=-1 \Leftrightarrow y=-1/x$. Comprobación en la ecuación original con $y=-1/x$, $y'=1/x^2$: $y+xy'=-\frac1x+\frac1x=0$ y $a(1+xy)=a(1-1)=0$, luego $y=-1/x$ es también solución, válida para \emph{cualquier} $a$.}
\begin{align*}
\int \frac{dz}{1+z} &= a\int dx \\
\ln|1+z| &= ax + C \implies 1+z = K e^{ax} \implies 1+xy = K e^{ax}
\end{align*}
Aplicando $y(1/a)=-a \implies 1 + \frac{1}{a}(-a) = K e^{a(1/a)} \implies 0 = K e \implies K=0$, valor que coincide con el hallado arriba.
\[
1+xy = 0 \implies
\]
\[
\boxed{y = -\frac{1}{x}}
\]


\noindent\textbf{12. Integrad la EDO: $\bm{y'+\sin(x-y)=\sin(x+y)}$ con la condición inicial $\bm{y(\pi)=\pi/2}$} \\
Desarrollamos la diferencia de senos paso a paso:
\begin{align*}
\sin(x+y) &= \sin(x)\cos(y) + \cos(x)\sin(y) \\
\sin(x-y) &= \sin(x)\cos(y) - \cos(x)\sin(y)
\end{align*}
\begin{align*}
\sin(x+y) - \sin(x-y) &= \big[\sin(x)\cos(y) + \cos(x)\sin(y)\big] - \big[\sin(x)\cos(y) - \cos(x)\sin(y)\big] \\
&= 2\cos(x)\sin(y)
\end{align*}
Por lo tanto:
\begin{align*}
y' &= 2\cos(x)\sin(y) \implies \frac{dy}{\sin(y)} = 2\cos(x)dx
\end{align*}
\textit{Al dividir por $\sin(y)$ se excluye $y=k\pi$; sustituyendo $y=k\pi$ (constante, $y'=0$) en la ecuación original: $\sin(x-k\pi)=\sin(x+k\pi)=\pm\sin(x)$ (mismo signo en ambos lados), por lo que $0+\sin(x-k\pi)=\sin(x+k\pi)$ se cumple, y $y=k\pi$ es también solución.}

Para integrar $\int \csc(y)\,dy$ : $u=\tan\left(\dfrac{y}{2}\right) \implies dy=\dfrac{2\,du}{1+u^2},\ \sin(y)=\dfrac{2u}{1+u^2}$.
\[
\csc(y)\,dy = \frac{1+u^2}{2u}\cdot\frac{2\,du}{1+u^2} = \frac{du}{u}
\]
\begin{align*}
\int \csc(y)\,dy &= \int \frac{du}{u} = \ln|u| = \ln\left|\tan\left(\frac{y}{2}\right)\right|
\end{align*}
\[
\ln\left|\tan\left(\frac{y}{2}\right)\right| = 2\sin(x) + C
\]
Aplicando $y(\pi)=\pi/2 \implies \ln|\tan(\pi/4)| = 2\sin(\pi)+C \implies \ln(1) = C \implies C=0$.
\[
\boxed{\tan\left(\frac{y}{2}\right) = e^{2\sin(x)}}
\]


\noindent\textbf{13. Integrad la EDO: $\bm{(1+x^2)y' - \frac{1}{2}\cos^2(2y) = 0}$ con $\bm{y \to 7\pi/2}$ cuando $\bm{x \to -\infty}$} \\
\begin{align*}
\frac{2\,dy}{\cos^2(2y)} &= \frac{dx}{1+x^2}
\end{align*}
\textit{Al dividir por $\cos^2(2y)$ se excluye $\cos(2y)=0$, es decir $y=\dfrac{\pi}{4}+\dfrac{k\pi}{2}$; en ese caso $y'=0$ y la ecuación original $(1+x^2)\cdot0-\frac12\cdot0=0$ se cumple, luego esas rectas son también solución.}
\begin{align*}
\int 2\sec^2(2y)\,dy &= \int \frac{dx}{1+x^2} \\
\tan(2y) &= \arctan(x) + C
\end{align*}
Límite $x \to -\infty \implies \arctan(x) \to -\pi/2$ e $y \to 7\pi/2$. \\
$\tan(7\pi) = -\frac{\pi}{2} + C \implies 0 = -\frac{\pi}{2} + C \implies C = \frac{\pi}{2}$.
\[
\boxed{\tan(2y) = \arctan(x) + \frac{\pi}{2}}
\]


\noindent\textbf{14. Integrad la EDO: $\bm{y' = 2x(\pi+y)}$ con $\bm{y}$ acotada cuando $\bm{x \to +\infty}$} \\
\begin{align*}
\frac{dy}{\pi+y} &= 2x\,dx
\end{align*}
\textit{Al dividir por $\pi+y$ se excluye $y=-\pi$; sustituyendo, $y'=0$ y $2x(\pi-\pi)=0$ se cumple para todo $x$, luego $y=-\pi$ es solución.}
\begin{align*}
\int \frac{dy}{\pi+y} &= \int 2x\,dx \\
\ln|\pi+y| &= x^2 + C \implies \pi+y = \pm K e^{x^2} \implies y = \pm K e^{x^2} - \pi
\end{align*}
Para que $y$ esté acotada cuando $x \to +\infty$, $K=0$
\[
\boxed{y = -\pi}
\]


\noindent\textbf{15. Integrad la EDO: $\bm{xy' = y + x\cos^2(y/x)}$} \\
$z = \dfrac{y}{x} \implies y = zx \implies y' = z'x+z$.
\begin{align*}
z'x + z &= z + \cos^2(z) \implies z'x = \cos^2(z) \implies \frac{dz}{\cos^2(z)} = \frac{dx}{x}
\end{align*}
\textit{La división por $\cos^2(z)$ excluye $\cos(z)=0$, es decir $z=\dfrac{\pi}{2}+k\pi$; en ese punto $z'x=0$ es consistente con $z$ constante, luego $y=x\left(\dfrac{\pi}{2}+k\pi\right)$ es solución.}
\begin{align*}
\int \sec^2(z)\,dz &= \int \frac{dx}{x} \\
\tan(z) &= \ln|x| + C
\end{align*}
\[
\boxed{\tan\left(\frac{y}{x}\right) = \ln|x| + C}
\]


\noindent\textbf{16. Integrad la EDO: $\bm{xy' = y(\ln y - \ln x)}$} \\
$z = \dfrac{y}{x} \implies y = zx \implies y' = z'x+z$; la ecuación se escribe $y' = z\ln(z)$.
\begin{align*}
z'x + z &= z\ln(z) \implies z'x = z(\ln(z)-1) \implies \frac{dz}{z(\ln(z)-1)} = \frac{dx}{x}
\end{align*}
\textit{El dominio exige $y>0$, es decir $z>0$. Al dividir por $\ln(z)-1$ se excluye $z=e$; en ese punto $z'x=e(1-1)=0$, consistente con $z$ constante, luego $y=ex$ es también solución.}

$u = \ln(z)-1 \implies du = \dfrac{dz}{z}$.
\begin{align*}
\int \frac{du}{u} &= \int \frac{dx}{x} \\
\ln|u| &= \ln|x| + C \implies \ln(z)-1 = Kx \implies \ln(y/x) = 1 + Kx
\end{align*}
$K=0$ da $y=xe^1=ex$, recuperando el caso anterior.
\[
\boxed{y = x\, e^{1+Kx}}
\]


\noindent\textbf{17. Integrad la EDO: $\bm{2x^2y' = x^2+y^2}$} \\
$z = \dfrac{y}{x} \implies y = zx \implies y' = z'x+z$.
\begin{align*}
z'x + z &= \frac{1+z^2}{2} \implies z'x = \frac{(z-1)^2}{2} \implies \frac{2\,dz}{(z-1)^2} = \frac{dx}{x}
\end{align*}
\textit{Al dividir por $(z-1)^2$ se excluye $z=1$; en ese punto $z'x=0$, consistente con $z$ constante. Así, $y=x$ es también solución, pero \emph{no} está contenida en la familia general (correspondería a $C\to\pm\infty$ en la expresión final).}
\begin{align*}
\int \frac{2\,dz}{(z-1)^2} &= \int \frac{dx}{x} \\
-\frac{2}{z-1} &= \ln|x| + C
\end{align*}
\[
\boxed{\frac{2x}{x-y} = \ln|x| + C}
\]


\noindent\textbf{18. Integrad la EDO: $\bm{(4x-3y)dx + (2y-3x)dy = 0}$} \\
$M=4x-3y$, $N=2y-3x$. $\dfrac{\partial M}{\partial y} = -3 = \dfrac{\partial N}{\partial x}$. Buscamos la función $f(x,y)$:
\begin{align*}
f(x,y) &= \int (4x-3y)\, dx = 2x^2 - 3xy + g(y) \\
\frac{\partial f}{\partial y} &= -3x + g'(y) = 2y - 3x \implies g'(y) = 2y \implies g(y) = y^2 +K
\end{align*}

\[
\boxed{2x^2 - 3xy + y^2 = C}
\]


\noindent\textbf{19. Integrad la EDO: $\bm{x+y-2 + (1-x)y' = 0}$} \\
Escribimos la ecuación en forma $(x+y-2)dx + (1-x)dy = 0$\\
$x=u+h$, $y=v+k$, y elegimos $h,k$ de modo que se anulen los términos independientes:
\[
\begin{cases} a_1 h + b_1 k + c_1 = 0 \\ a_2 h + b_2 k + c_2 = 0 \end{cases}
\implies
\begin{cases} h + k - 2 = 0 \\ -h + 1 = 0 \end{cases}
\]
Determinante del sistema:
\[
D = \begin{vmatrix} a_1 & b_1 \\ a_2 & b_2 \end{vmatrix} = \begin{vmatrix} 1 & 1 \\ -1 & 0 \end{vmatrix} = 1
\]
Como $D\neq 0$, el sistema tiene solución única. $h=k=1$.

\[
\frac{dv}{du} = \frac{u+v}{u} = 1+\frac{v}{u}
\]
$z = \dfrac{v}{u} \implies v=zu \implies \dfrac{dv}{du}=z'u+z$.
\begin{align*}
z'u+z &= 1+z \implies z'u = 1 \implies \int dz = \int \frac{du}{u} \implies z = \ln|u| + C
\end{align*}
\[
\frac{y-1}{x-1} = \ln|x-1| + C
\]
\[
\boxed{y = 1 + (x-1)\big(\ln|x-1| + C\big)}
\]


\noindent\textbf{20. Integrad la EDO: $\bm{(3y-7x+7)dx - (3x-7y-3)dy = 0}$} \\
Analogamente al 19. :
\[
\begin{cases} a_1 h + b_1 k + c_1 = 0 \\ a_2 h + b_2 k + c_2 = 0 \end{cases}
\implies
\begin{cases} -7h + 3k = -7 \\ -3h + 7k = -3 \end{cases}
\]
Determinante del sistema:
\[
D = \begin{vmatrix} a_1 & b_1 \\ a_2 & b_2 \end{vmatrix} = \begin{vmatrix} -7 & 3 \\ -3 & 7 \end{vmatrix} = -49+9 = -40
\]

\[
-49h+21k=-49, \qquad -9h+21k=-9 \implies -40h=-40 \implies h=1
\]
\[
h=1, \qquad k=0
\]

\[
\frac{dv}{du} = \frac{-7u+3v}{3u-7v} = \frac{-7+3(v/u)}{3-7(v/u)}
\]
$z = \dfrac{v}{u} \implies v=zu \implies \dfrac{dv}{du}=z'u+z$.
\begin{align*}
z'u+z &= \frac{3z-7}{3-7z} \implies z'u = \frac{3z-7-z(3-7z)}{3-7z} = \frac{7(z^2-1)}{3-7z}
\end{align*}
\textit{Al dividir por $u=x-1$ se excluye $x=1$, correspondiente al punto singular $(x,y)=(1,0)$ (el centro del sistema). Además, la división por $(3-7z)$ excluye $z=3/7$ (allí las curvas tienen tangente vertical, no es una solución). Por otro lado, cuando $z^2-1=0$ (es decir $z=\pm1$) se tiene $z'u=0$, consistente con $z$ constante:}
\begin{itemize}
\item $z=1 \Rightarrow v=u \Rightarrow y=x-1$. Comprobación: con $dy=dx$, $(3y-7x+7)-( 3x-7y-3)=(-4x+4)-(-4x+4)=0$ \checkmark.
\item $z=-1 \Rightarrow v=-u \Rightarrow y=1-x$. Comprobación: con $dy=-dx$, $(3y-7x+7)+(3x-7y-3)=-4y-4x+4=0$ \checkmark.
\end{itemize}
Ambas rectas son, por tanto, soluciones singulares.

\[
\frac{3-7z}{7(z^2-1)}dz = (\dfrac{-2/7}{z-1} - \dfrac{5/7}{z+1} )dz = \frac{du}{u}
\]

\begin{align*}
\int \left(\frac{-2/7}{z-1} - \frac{5/7}{z+1}\right)dz &= \int \frac{du}{u} \\
-\frac{2}{7}\ln|z-1| - \frac{5}{7}\ln|z+1| &= \ln|u| + K_1 \\
7\ln|u| + 2\ln|z-1| + 5\ln|z+1| &= C
\end{align*}

\[
u^7\left(\frac{v}{u}-1\right)^2\left(\frac{v}{u}+1\right)^5 = C_2 \implies (v-u)^2(v+u)^5 = C_2
\]
\[
\boxed{(y-x+1)^2(y+x-1)^5 = C_3}
\]

\separador
\section{Hoja 2}

\noindent\textbf{1. Integrad la EDO: $\bm{8x+4y+1+(4x+2y+1)y^{\prime}=0}$} \\
$z = 4x+2y \implies z' = 4+2y' \implies y' = \dfrac{z'-4}{2}$.
\begin{align*}
2z+1 + (z+1)\frac{z'-4}{2} &= 0 \implies 2(2z+1) + (z+1)(z'-4) = 0 \\
4z+2 + (z+1)z' - 4z - 4 &= 0 \implies (z+1)z' = 2 \implies (z+1)dz = 2dx
\end{align*}
\begin{align*}
\int (z+1)dz &= \int 2dx \implies \frac{1}{2}z^2 + z = 2x + C_1 \implies z^2 + 2z = 4x + C
\end{align*}
\begin{align*}
(4x+2y)^2 + 2(4x+2y) &= 4x + C \implies 16x^2 + 16xy + 4y^2 + 8x + 4y = 4x + C
\end{align*}
\[
\boxed{(4x+2y)^2 + 4x + 4y = C}
\]


\noindent\textbf{2. Integrad la EDO: $\bm{(x-2y-1)dx+(3x-6y+2)dy=0}$} \\
$z = x-2y \implies dx = dz + 2dy$.
\begin{align*}
(z-1)(dz + 2dy) + (3z+2)dy &= 0 \implies (z-1)dz + (2z-2+3z+2)dy = 0 \\
(z-1)dz + 5z\,dy &= 0 \implies \left(1 - \frac{1}{z}\right)dz = -5dy
\end{align*}
\textit{Al dividir por $z$ se excluye $z=0$, es decir $x-2y=0 \implies x=2y \implies dx=2dy$. Sustituyendo $(-1)(2dy) + 2dy = 0$, luego $x-2y=0$ es solución.}
\begin{align*}
\int \left(1 - \frac{1}{z}\right)dz &= \int -5dy \implies z - \ln|z| = -5y + C_1 \\
(x-2y) - \ln|x-2y| &= -5y + C_1 \implies x + 3y - \ln|x-2y| = C_1
\end{align*}
$\implies x-2y = \pm K e^{x+3y}$ con $K>0$; $K\geq 0$ incluye la $x=2y$.
\[
\boxed{x - 2y = \pm K e^{x+3y}, \qquad K \geq 0}
\]


\noindent\textbf{3. Integrad la EDO: $\bm{y^{\prime}+2y=e^{-x}}$} \\
\begin{align*}
y'_h + 2y_h &= 0 \implies \frac{dy_h}{y_h} = -2\,dx \implies \ln|y_h| = -2x + C_1 \implies y_h(x) = K e^{-2x}
\end{align*}
$y(x) = C(x) e^{-2x} \implies y' = C'(x) e^{-2x} - 2C(x) e^{-2x}$. \\
\begin{align*}
C'(x) e^{-2x} - 2C(x) e^{-2x} + 2C(x) e^{-2x} &= e^{-x} \implies C'(x) e^{-2x} = e^{-x} \implies C'(x) = e^x \\
C(x) &= \int e^x dx = e^x + C
\end{align*}
\[
\boxed{y(x) = e^{-x} + C e^{-2x}}
\]


\noindent\textbf{4. Integrad la EDO: $\bm{x^{2}+xy^{\prime}=y}$ con la condición inicial $\bm{y(x=1)=0}$} \\
$y' - \frac{1}{x}y = -x$. \\
\begin{align*}
y'_h - \frac{1}{x}y_h &= 0 \implies \frac{dy_h}{y_h} = \frac{dx}{x} \implies \ln|y_h| = \ln|x| + C_1 \implies y_h(x) = Kx
\end{align*}
$y(x) = C(x)x \implies y' = C'(x)x + C(x)$. \\
\begin{align*}
C'(x)x + C(x) - \frac{1}{x}C(x)x &= -x \implies C'(x)x = -x \implies C'(x) = -1 \\
C(x) &= -x + C_0 \implies y(x) = (C_0 - x)x = C_0 x - x^2
\end{align*}
\textit{Al dividir por $x$ se excluye $x=0$} \\
$y(1)=0 \implies C_0(1) - 1^2 = 0 \implies C_0 = 1$.
\[
\boxed{y(x) = x - x^2}
\]


\noindent\textbf{5. Integrad la EDO: $\bm{y^{\prime}\cos x-y\sin x=2x}$ con la condición inicial $\bm{y(x=0)=0}$} \\
Dividiendo por $\cos x$: $y' - y\tan x = \frac{2x}{\cos x}$. \\
\begin{align*}
y'_h - y_h\tan x &= 0 \implies \frac{dy_h}{y_h} = \tan x\,dx \implies \ln|y_h| = -\ln|\cos x| + C_1 \implies y_h(x) = K\sec x
\end{align*}
$y(x) = C(x)\sec x \implies y' = C'(x)\sec x + C(x)\sec x\tan x$. \\
Sustituyendo en $y'\cos x - y\sin x = 2x$:
\begin{align*}
\big(C'(x)\sec x + C(x)\sec x\tan x\big)\cos x - C(x)\sec x\sin x &= 2x \implies C'(x) = 2x \\
C(x) &= x^2 + C_0 \implies y(x) = (x^2 + C_0)\sec x
\end{align*}
\textit{La EDO presenta singularidades en $x = \frac{\pi}{2} + k\pi$ donde $\cos x = 0$.} \\
$y(0)=0 \implies (0^2 + C_0)\sec(0) = 0 \implies C_0 = 0$.
\[
\boxed{y(x) = x^2 \sec x}
\]


\noindent\textbf{6. Integrad la EDO: $\bm{xy^{\prime}-2y=x^{3}\cos x}$} \\
$y' - \frac{2}{x}y = x^2 \cos x$. \\
\begin{align*}
y'_h - \frac{2}{x}y_h &= 0 \implies \frac{dy_h}{y_h} = \frac{2}{x}dx \implies \ln|y_h| = 2\ln|x| + C_1 \implies y_h(x) = Kx^2
\end{align*}
$y(x) = C(x)x^2 \implies y' = C'(x)x^2 + 2C(x)x$. \\
\begin{align*}
x\big(C'(x)x^2 + 2C(x)x\big) - 2C(x)x^2 &= x^3 \cos x \implies C'(x)x^3 = x^3 \cos x \implies C'(x) = \cos x \\
C(x) &= \sin x + C_0 \implies y(x) = (\sin x + C_0)x^2
\end{align*}
\textit{$y=0$ es solución.}
\[
\boxed{y(x) = x^2 \sin x + C x^2}
\]


\noindent\textbf{7. Integrad la EDO: $\bm{(2x-y^{2})y^{\prime}=2y}$} \\
$P(x,y)dx + N(x,y)dy = 0$:
\begin{align*}
-2y\,dx + (2x-y^2)dy &= 0 \implies P(x,y) = -2y, \quad N(x,y) = 2x-y^2
\end{align*}
\begin{align*}
\frac{\partial P}{\partial y} = -2 \neq \frac{\partial N}{\partial x} = 2 \implies \text{No es exacta.}
\end{align*}
$\mu(y)$:
\begin{align*}
\frac{\frac{\partial N}{\partial x} - \frac{\partial P}{\partial y}}{P} = \frac{2 - (-2)}{-2y} = -\frac{2}{y} \implies \mu(y) = e^{\int -\frac{2}{y}dy} = e^{-2\ln|y|} = \frac{1}{y^2}
\end{align*}
\begin{align*}
-\frac{2}{y}dx + \left(\frac{2x}{y^2} - 1\right)dy &= 0
\end{align*}
Comprobación enla nueva EDO:
\begin{align*}
\tilde{P}(x,y) = -\frac{2}{y} \implies \frac{\partial \tilde{P}}{\partial y} = \frac{2}{y^2}, \qquad \tilde{N}(x,y) = \frac{2x}{y^2} - 1 \implies \frac{\partial \tilde{N}}{\partial x} = \frac{2}{y^2} \quad
\end{align*}
\begin{align*}
F(x,y) &= \int \tilde{P}(x,y) dx = \int -\frac{2}{y}dx = -\frac{2x}{y} + g(y) \\
\frac{\partial F}{\partial y} &= \frac{2x}{y^2} + g'(y) = \tilde{N}(x,y) = \frac{2x}{y^2} - 1 \implies g'(y) = -1 \implies g(y) = -y
\end{align*}
La solución implícita $F(x,y) = C_1$ es:
\begin{align*}
-\frac{2x}{y} - y = C_1 \implies 2x + y^2 = Ky
\end{align*}
\textit{Al dividir por $y^2$ para el factor integrante se excluye $y=0$. Sustituyendo $y\equiv 0$ ($y'=0$) en la EDO original: $(2x-0)\cdot 0 = 2(0)$ se cumple $\forall x$, luego $y=0$ es solución, recuperada con $K=0$.}
\[
\boxed{2x + y^2 = Ky}
\]

\noindent\textbf{8. Integrad la EDO: $\bm{y^{\prime}=\frac{y}{2y\ln y+y-x}}$} \\
$\frac{dx}{dy} + \frac{1}{y}x = 2\ln y + 1$. \\

\begin{align*}
\frac{dx_h}{dy} + \frac{1}{y}x_h &= 0 \implies \frac{dx_h}{x_h} = -\frac{dy}{y} \implies x_h(y) = \frac{C_1}{y}
\end{align*}
$x(y) = \frac{C(y)}{y} \implies \frac{dx}{dy} = \frac{C'(y)y - C(y)}{y^2}$. \\
\begin{align*}
\frac{C'(y)y - C(y)}{y^2} + \frac{C(y)}{y^2} &= 2\ln y + 1 \implies \frac{C'(y)}{y} = 2\ln y + 1 \implies C'(y) = 2y\ln y + y \\
C(y) &= \int (2y\ln y + y)dy = y^2\ln y - \frac{y^2}{2} + \frac{y^2}{2} + K = y^2\ln y + K
\end{align*}
\textit{El dominio pide $y > 0$}
\[
\boxed{x(y) = y\ln y + \frac{K}{y}}
\]


\noindent\textbf{9. Integrad la EDO: $\bm{(e^{-y^{2}/2}-xy)dy-dx=0}$} \\
$\frac{dx}{dy} + yx = e^{-y^2/2}$. \\
\begin{align*}
\frac{dx_h}{dy} + yx_h &= 0 \implies \frac{dx_h}{x_h} = -y\,dy \implies x_h(y) = C_1 e^{-y^2/2}
\end{align*}
$x(y) = C(y)e^{-y^2/2} \implies \frac{dx}{dy} = C'(y)e^{-y^2/2} - C(y)y e^{-y^2/2}$. \\
\begin{align*}
C'(y)e^{-y^2/2} - C(y)y e^{-y^2/2} + y C(y)e^{-y^2/2} &= e^{-y^2/2} \implies C'(y) = 1 \implies C(y) = y + C
\end{align*}
\[
\boxed{x e^{y^2/2} - y = C}
\]


\noindent\textbf{10. Integrad la EDO: $\bm{y^{\prime}-y\ln 2=2^{\sin x}(\cos x-1)\ln 2}$ con $\bm{y}$ acotada cuando $\bm{x\rightarrow+\infty}$} \\
\begin{align*}
y'_h - y_h\ln 2 &= 0 \implies \frac{dy_h}{y_h} = \ln 2\,dx \implies y_h(x) = K \cdot 2^x
\end{align*}
$y(x) = C(x) \cdot 2^x \implies y' = C'(x)2^x + C(x)2^x \ln 2$. \\
\begin{align*}
C'(x)2^x + C(x)2^x \ln 2 - C(x)2^x \ln 2 &= 2^{\sin x}(\cos x - 1)\ln 2 \implies C'(x) = 2^{\sin x - x}(\cos x - 1)\ln 2
\end{align*}
$u = \sin x - x \implies du = (\cos x - 1)dx$:
\begin{align*}
C(x) &= \int 2^u \ln 2\,du = 2^u + C_0 = 2^{\sin x - x} + C_0 \implies y(x) = 2^{\sin x} + C_0 \cdot 2^x
\end{align*}

\textit{El término $2^{\sin x}$ está acotado en $[1/2, 2]\ \forall x \in \mathbb{R}$. Analizando la acotación cuando $x \to +\infty$:}
\begin{itemize}
\item \textit{Si $C_0 > 0$, la función cumple acotación por abajo ($y(x) \geq 1/2$), pero diverge a $+\infty$ por arriba.}
\item \textit{Si $C_0 < 0$, cumple acotación superior ($y(x) \leq 2$), pero diverge a $-\infty$ por abajo.}
\item \textit{Por tanto, para asegurar acotación global (tanto superior como inferior) se requiere $C_0 = 0$.}
\end{itemize}
\[
\boxed{y(x) = 2^{\sin x}}
\]


\noindent\textbf{11. Integrad la EDO: $\bm{y^{\prime}+2xy=2xy^{2}}$} \\
$z = y^{-1} \implies z' = -y^{-2}y' \implies z' - 2xz = -2x$. \\
\begin{align*}
z'_h - 2xz_h &= 0 \implies \frac{dz_h}{z_h} = 2x\,dx \implies z_h(x) = K e^{x^2}
\end{align*}
$z(x) = C(x)e^{x^2} \implies z' = C'(x)e^{x^2} + 2xC(x)e^{x^2}$. \\
\begin{align*}
C'(x)e^{x^2} + 2xC(x)e^{x^2} - 2xC(x)e^{x^2} &= -2x \implies C'(x) = -2x e^{-x^2} \\
C(x) &= e^{-x^2} + C_0 \implies z(x) = (e^{-x^2} + C_0)e^{x^2} = 1 + C_0 e^{x^2}
\end{align*}
\textit{Al dividir por $y^2$ se excluye $y=0$. $y\equiv 0$ en la EDO original, $0' + 2x(0) = 2x(0)^2$ se satisface $\forall x$, luego $y=0$ es solución (alcanzada para $C_0 \to \infty$).}
\[
\boxed{y(x) = \frac{1}{1 + C e^{x^2}}}
\]


\noindent\textbf{12. Integrad la EDO: $\bm{(x^{3}+e^{y})y^{\prime}=3x^{2}}$} \\
$3x^2 \frac{dx}{dy} - x^3 = e^y$. \\
$z = x^3 \implies \frac{dz}{dy} = 3x^2 \frac{dx}{dy} \implies \frac{dz}{dy} - z = e^y$. \\
\begin{align*}
\frac{dz_h}{dy} - z_h &= 0 \implies \frac{dz_h}{z_h} = dy \implies z_h(y) = K e^y
\end{align*}
$z(y) = C(y)e^y \implies \frac{dz}{dy} = C'(y)e^y + C(y)e^y$. \\
\begin{align*}
C'(y)e^y + C(y)e^y - C(y)e^y &= e^y \implies C'(y) = 1 \implies C(y) = y + C_0 \\
z(y) &= (y + C_0)e^y \implies x^3 = (y + C_0)e^y
\end{align*}
\[
\boxed{x^3 = (y+C)e^y}
\]

\noindent\textbf{13. Integrad la EDO: $\bm{y^{\prime}-2ye^{x}=2\sqrt{y}e^{x}}$} \\
$n=1/2$. $z = \sqrt{y} \implies z' = \frac{1}{2}y^{-1/2}y' \implies z' - z e^x = e^x$. \\
\begin{align*}
z'_h - z_h e^x &= 0 \implies \frac{dz_h}{z_h} = e^x dx \implies z_h(x) = K e^{e^x}
\end{align*}
$z(x) = C(x)e^{e^x} \implies z' = C'(x)e^{e^x} + C(x)e^x e^{e^x}$. \\
\begin{align*}
C'(x)e^{e^x} + C(x)e^x e^{e^x} - C(x)e^x e^{e^x} &= e^x \implies C'(x) = e^x e^{-e^x} \\
C(x) &= -e^{-e^x} + C_0 \implies z(x) = (C_0 - e^{-e^x})e^{e^x} = C_0 e^{e^x} - 1
\end{align*}
\textit{Al dividir por $\sqrt{y}$ se excluye $y=0$; $y\equiv 0$ es solución. Como $z = \sqrt{y} \geq 0$, exige la restricción $C e^{e^x} \geq 1$.}
\[
\boxed{y(x) = \left(C e^{e^x} - 1\right)^2, \qquad C e^{e^x} \geq 1}
\]


\noindent\textbf{14. Integrad la EDO: $\bm{2y^{\prime}\ln x+\frac{y}{x}=\frac{\cos x}{y}}$} \\
$2y y' \ln x + \frac{y^2}{x} = \cos x$. \\
$z = y^2 \implies z' = 2y y' \implies z' \ln x + \frac{z}{x} = \cos x$. \\
\begin{align*}
z'_h \ln x + \frac{z_h}{x} &= 0 \implies \frac{dz_h}{z_h} = -\frac{dx}{x\ln x} \implies \ln|z_h| = -\ln|\ln x| + C_1 \implies z_h(x) = \frac{K}{\ln x}
\end{align*}
Proponemos $z(x) = \frac{C(x)}{\ln x} \implies z' = \frac{C'(x)\ln x - C(x)/x}{\ln^2 x}$. \\
Sustituyendo en la EDO reducida:
\begin{align*}
\left(\frac{C'(x)\ln x - C(x)/x}{\ln^2 x}\right)\ln x + \frac{C(x)}{x\ln x} &= \cos x \implies \frac{C'(x)}{\ln x} = \cos x \implies C'(x) = \cos x \\
C(x) &= \sin x + C_0 \implies z(x) = \frac{\sin x + C_0}{\ln x}
\end{align*}
\textit{El dominio requiere $x > 0$, $x \neq 1$. La división por $y$ excluye $y=0$.}
\[
\boxed{y^2 = \frac{\sin x + C}{\ln x}}
\]


\noindent\textbf{15. Integrad la EDO: $\bm{2y^{\prime}\sin x+y\cos x=y^{3}\sin^{2}x}$} \\
Bernoulli con $n=3$. $2y^{-3}y'\sin x + y^{-2}\cos x = \sin^2 x$. \\
$z = y^{-2} \implies z' = -2y^{-3}y' \implies -z'\sin x + z\cos x = \sin^2 x$. \\
\begin{align*}
-z'_h \sin x + z_h \cos x &= 0 \implies \frac{dz_h}{z_h} = \cot x\,dx \implies z_h(x) = K \sin x
\end{align*}
$z(x) = C(x)\sin x \implies z' = C'(x)\sin x + C(x)\cos x$. \\
\begin{align*}
-\big(C'(x)\sin x + C(x)\cos x\big)\sin x + C(x)\sin x\cos x &= \sin^2 x \implies -C'(x)\sin^2 x = \sin^2 x \\
C'(x) &= -1 \implies C(x) = C_0 - x \implies z(x) = (C_0 - x)\sin x
\end{align*}
\textit{Al dividir por $y^3$ se excluye $y=0$, que es solución }
\[
\boxed{y^2 = \frac{1}{(C - x)\sin x}}
\]


\noindent\textbf{16. Resolved mediante derivación la ecuación: $\bm{\int_{0}^{x}ty(t)dt=x^{2}y(x)}$} \\
\begin{align*}
x y(x) &= 2x y(x) + x^2 y'(x) \implies x^2 y'(x) + x y(x) = 0
\end{align*}
Con $x \neq 0$:
\begin{align*}
x y' + y = 0 \implies \frac{dy}{y} = -\frac{dx}{x} \implies \ln|y| = -\ln|x| + C_1
\end{align*}
\[
\boxed{y(x) = \pm\frac{C}{x}, \qquad C>0}
\]


\noindent\textbf{17. Resolved mediante derivación la ecuación: $\bm{\int_{0}^{1}y(xt)dt=ny(x)}$} \\
$u = xt \implies du = x\,dt \implies dt = \frac{du}{x}$.\\
$t=0 \to u=0$, $t=1 \to u=x$.
\begin{align*}
\frac{1}{x} \int_{0}^{x} y(u) du &= n y(x) \implies \int_{0}^{x} y(u) du = n x y(x)
\end{align*}
\begin{align*}
y(x) &= n \left[ y(x) + x y'(x) \right] \implies y(x) = n y(x) + n x y'(x) \\
n x y'(x) &= (1-n) y(x) \implies \frac{dy}{y} = \frac{1-n}{n} \frac{dx}{x}
\end{align*}
$\ln|y| = \frac{1-n}{n}\ln|x| + C_1$.
\[
\boxed{y(x) = C x^{\frac{1-n}{n}}}
\]


\noindent\textbf{18. Integrad la EDO: $\bm{x(2x^{2}+y^{2})+y(x^{2}+2y^{2})y^{\prime}=0}$} \\
$M(x,y)dx + N(x,y)dy = 0$:
\begin{align*}
(2x^3 + xy^2)dx + (x^2y + 2y^3)dy &= 0
\end{align*}
Comprobación de exacta: $\frac{\partial M}{\partial y} = 2xy = \frac{\partial N}{\partial x} \implies \text{Lo es, xd}$. \\
\begin{align*}
f(x,y) &= \int (2x^3 + xy^2)dx = \frac{1}{2}x^4 + \frac{1}{2}x^2y^2 + g(y) \\
\frac{\partial f}{\partial y} &= x^2y + g'(y) = x^2y + 2y^3 \implies g'(y) = 2y^3 \implies g(y) = \frac{1}{2}y^4
\end{align*}
\[
\boxed{x^4 + x^2y^2 + y^4 = C}
\]


\noindent\textbf{19. Integrad la EDO: $\bm{\left(\frac{x}{\sqrt{x^{2}+y^{2}}}+\frac{1}{x}+\frac{1}{y}\right)dx+\left(\frac{y}{\sqrt{x^{2}+y^{2}}}+\frac{1}{y}-\frac{x}{y^{2}}\right)dy=0}$} \\
\begin{align*}
\frac{\partial M}{\partial y} &= -\frac{xy}{(x^2+y^2)^{3/2}} - \frac{1}{y^2}, \qquad \frac{\partial N}{\partial x} = -\frac{xy}{(x^2+y^2)^{3/2}} - \frac{1}{y^2} \implies \text{Es exacta}
\end{align*}
\begin{align*}
f(x,y) &= \int \left(\frac{x}{\sqrt{x^2+y^2}} + \frac{1}{x} + \frac{1}{y}\right) dx = \sqrt{x^2+y^2} + \ln|x| + \frac{x}{y} + g(y) \\
\frac{\partial f}{\partial y} &= \frac{y}{\sqrt{x^2+y^2}} - \frac{x}{y^2} + g'(y) = \frac{y}{\sqrt{x^2+y^2}} + \frac{1}{y} - \frac{x}{y^2} \implies g'(y) = \frac{1}{y} \implies g(y) = \ln|y|
\end{align*}
\textit{El dominio de la EDO requiere $x \neq 0$ y $y \neq 0$.}
\[
\boxed{\sqrt{x^2+y^2} + \ln|xy| + \frac{x}{y} = C}
\]


\noindent\textbf{20. Integrad la EDO: $\bm{\left(3x^{2}\tan y-\frac{2y^{3}}{x^{3}}\right)dx+\left(x^{3}\sec^{2}y+4y^{3}+\frac{3y^{2}}{x^{2}}\right)dy=0}$} \\
\begin{align*}
\frac{\partial M}{\partial y} &= 3x^2\sec^2 y - \frac{6y^2}{x^3}, \qquad \frac{\partial N}{\partial x} = 3x^2\sec^2 y - \frac{6y^2}{x^3} \implies \text{Otra vez exacta}\\
f(x,y) &= \int \left(3x^2\tan y - \frac{2y^3}{x^3}\right) dx = x^3\tan y + \frac{y^3}{x^2} + g(y) \\
\frac{\partial f}{\partial y} &= x^3\sec^2 y + \frac{3y^2}{x^2} + g'(y) = x^3\sec^2 y + 4y^3 + \frac{3y^2}{x^2} \implies g'(y) = 4y^3 \implies g(y) = y^4
\end{align*}
\textit{El dominio pide $x \neq 0$ y $y \neq \frac{\pi}{2} + k\pi$.}
\[
\boxed{x^3\tan y + \frac{y^3}{x^2} + y^4 = C}
\]

\separador
\section{Hoja 3}

\noindent\textbf{1. Integrad la EDO: $\bm{(x^4\ln x - 2xy^3)dx + 3x^2y^2dy = 0}$} \\
$M(x,y) = x^4\ln x - 2xy^3$ y $N(x,y) = 3x^2y^2$:
\begin{align*}
\frac{\partial M}{\partial y} = -6xy^2, \qquad \frac{\partial N}{\partial x} = 6xy^2 \implies \text{No es exacta.}
\end{align*}
\begin{align*}
\frac{\partial}{\partial y}\big(M(x,y)\mu(x)\big) = \frac{\partial}{\partial x}\big(N(x,y)\mu(x)\big)
\end{align*}
\begin{align*}
\mu(x)\frac{\partial M}{\partial y} &= \mu'(x)N(x,y) + \mu(x)\frac{\partial N}{\partial x}
\end{align*}
\begin{align*}
\mu(x)(-6xy^2) &= \mu'(x)(3x^2y^2) + \mu(x)(6xy^2)
\end{align*}
\begin{align*}
\mu'(x) &= \frac{-12xy^2}{3x^2y^2} \, \mu(x) \implies \mu'(x) = -\frac{4}{x} \, \mu(x)\implies \mu(x) = e^{\int -\frac{4}{x}dx} = \frac{1}{x^4}
\end{align*}

Comprobamos la exactitud de la nueva EDO:
\begin{align*}
\frac{\partial}{\partial y}\big(M(x,y)\mu(x)\big) &= \frac{\partial}{\partial y}\left(\ln x - \frac{2y^3}{x^3}\right) = -\frac{6y^2}{x^3} \\
\frac{\partial}{\partial x}\big(N(x,y)\mu(x)\big) &= \frac{\partial}{\partial x}\left(\frac{3y^2}{x^2}\right) = -\frac{6y^2}{x^3}
\end{align*}
Integrando $\tilde{N} = \frac{3y^2}{x^2}$ respecto de $y$ obtenemos $f(x,y) = \frac{y^3}{x^2} + h(x)$. Derivando e igualando a $\tilde{M}$:
\begin{align*}
\frac{\partial f}{\partial x} = -\frac{2y^3}{x^3} + h'(x) = \ln x - \frac{2y^3}{x^3} \implies h'(x) = \ln x \implies h(x) = x(\ln x - 1)
\end{align*}
\[
\boxed{\frac{y^3}{x^2} + x\ln x - x = C}
\]

\noindent\textbf{2. Integrad la EDO: $\bm{(x+\sin x+\sin y)dx + \cos y \, dy = 0}$} \\
$M(x,y) = x+\sin x+\sin y$ y $N(x,y) = \cos y$:
\begin{align*}
\frac{\partial M}{\partial y} = \cos y, \qquad \frac{\partial N}{\partial x} = 0 \implies \text{No es exacta.}
\end{align*}
\begin{align*}
\frac{\partial}{\partial y}\big(M(x,y)\mu(x)\big) = \frac{\partial}{\partial x}\big(N(x,y)\mu(x)\big)
\end{align*}
\begin{align*}
\mu(x)\cos y &= \mu'(x)\cos y + \mu(x)(0)
\end{align*}
\begin{align*}
\mu'(x) = \mu(x) \implies \mu(x) = e^x
\end{align*}

Comprobamos la exactitud de la nueva EDO:
\begin{align*}
\left.\begin{aligned}
\frac{\partial \tilde{M}}{\partial y} &= \frac{\partial}{\partial y}\big(e^x(x+\sin x+\sin y)\big) = e^x\cos y \\[1.5ex]
\frac{\partial \tilde{N}}{\partial x} &= \frac{\partial}{\partial x}(e^x\cos y) = e^x\cos y
\end{aligned}\right\} \implies \text{Oh sorpresa, es exacta.}
\end{align*}
Integrando $\tilde{N} = e^x\cos y$ respecto de $y$ tenemos $f(x,y) = e^x\sin y + h(x)$. Derivando e igualando a $\tilde{M}$:
\begin{align*}
\frac{\partial f}{\partial x} = e^x\sin y + h'(x) = e^x(x+\sin x+\sin y) \implies h'(x) = e^x(x+\sin x)
\end{align*}
Integrando $h'(x)$ por partes:
\begin{align*}
h(x) = \int e^x(x+\sin x) dx = e^x(x-1) + \frac{e^x(\sin x - \cos x)}{2}
\end{align*}
\[
\boxed{\sin y + x - 1 + \frac{\sin x - \cos x}{2} = C e^{-x}}
\]


\noindent\textbf{3. Integrad la EDO: $\bm{y = (y')^2 e^{y'}}$} \\
$y' = p \implies y = p^2 e^p$. Diferenciando $y$:
\begin{align*}
\frac{dy}{dp} &= (2pe^p + p^2 e^p)dp = p(p+2)e^p
\end{align*}
Como $dy = p dx$, ($p\neq 0$):
\begin{align*}
p dx &= p(p+2)e^p dp \implies dx = (p+2)e^p dp \\
x &= \int (p+2)e^p dp = (p+1)e^p + C
\end{align*}
\textit{Para $p=0$ se obtiene la solución $y=0$ (recuperada en el límite).}
\[
\boxed{\begin{cases} x = (p+1)e^p + C \\ y = p^2 e^p \end{cases}}
\]


\noindent\textbf{4. Integrad la EDO: $\bm{x = \ln y' + \sin y'}$} \\
$y' = p \implies x = \ln p + \sin p$
\begin{align*}
dx &= \left(\frac{1}{p} + \cos p\right)dp
\end{align*}
haciendo $dy = p dx$:
\begin{align*}
dy &= p\left(\frac{1}{p} + \cos p\right)dp = (1 + p\cos p)dp \\
y &= \int (1 + p\cos p)dp = p + p\sin p + \cos p + C
\end{align*}
\[
\boxed{\begin{cases} x = \ln p + \sin p \\ y = p + p\sin p + \cos p + C \end{cases}}
\]


\noindent\textbf{5. Integrad la EDO: $\bm{x(1+(y')^2)^{3/2} = a}$} \\
$y' = p \implies x = \dfrac{a}{(1+p^2)^{3/2}}$. Diferenciando $x$:
\begin{align*}
dx &= -\frac{3ap}{(1+p^2)^{5/2}}dp
\end{align*}
Como $dy = p dx$:
\begin{align*}
dy &= -\frac{3ap^2}{(1+p^2)^{5/2}}dp \implies y = -3a\int \frac{p^2}{(1+p^2)^{5/2}}dp
\end{align*}
$p = \tan\theta \implies dp = \sec^2\theta \, d\theta$, con lo que $x = a\cos^3\theta$

\begin{align*}
y &= -3a\int \sin^2\theta\cos\theta \, d\theta = C - a\sin^3\theta
\end{align*}
\[
\boxed{\begin{cases} x = a\cos^3\theta \\[1ex] y = C - a\sin^3\theta \end{cases}}
\]

\noindent\textbf{6. Integrad la EDO: $\bm{x = y' + \sin y'}$} \\
$y' = p \implies x = p + \sin p$.
\begin{align*}
dx &= (1 + \cos p)dp \quad\quad dy = p dx
\end{align*}

\begin{align*}
dy &= p(1 + \cos p)dp = (p + p\cos p)dp \\
y &= \int (p + p\cos p)dp = \frac{p^2}{2} + p\sin p + \cos p + C
\end{align*}
\[
\boxed{\begin{cases} x = p + \sin p \\ y = \dfrac{p^2}{2} + p\sin p + \cos p + C \end{cases}}
\]


\noindent\textbf{7. Integrad la EDO: $\bm{y = 2xy' + \ln y'}$} \\
$y' = p \implies y = 2xp + \ln p$.
\begin{align*}
dy &= 2p dx + 2xdp + \frac{1}{p}dp
\end{align*}
Sustituyendo $dy = pdx$:
\begin{align*}
p dx &= 2p dx + \left(2x + \frac{1}{p}\right)dp \implies \frac{dx}{dp} + \frac{2}{p}x = -\frac{1}{p^2}
\end{align*}
\textit{$x_h = \dfrac{K}{p^2}$ y variando la constante $x(p) = \dfrac{C-p}{p^2}$. Sustituyendo en la expresión de $y(p)$:}
\[
\boxed{\begin{cases} x = \dfrac{C-p}{p^2} \\[2ex] y = \dfrac{2(C-p)}{p} + \ln p \end{cases}}
\]


\noindent\textbf{8. Integrad la EDO: $\bm{y = 2xy' + \sin y'}$} \\
$y' = p \implies y = 2xp + \sin p$.
\begin{align*}
dy &= 2p dx + 2xdp + \cos p \, dp \qquad dy = pdx 
\end{align*}
\begin{align*}
p dx &= 2p dx + (2x + \cos p)dp \implies \frac{dx}{dp} + \frac{2}{p}x = -\frac{\cos p}{p}
\end{align*}
\textit{Resolvemos la lineal, $x_h = \dfrac{K(p)}{p^2}$ y $x(p) = \dfrac{C - p\sin p - \cos p}{p^2}$.}
\[
\boxed{\begin{cases} x = \dfrac{C - p\sin p - \cos p}{p^2} \\[2ex] y = \dfrac{2(C - \cos p)}{p} - \sin p \end{cases}}
\]


\noindent\textbf{9. Integrad la EDO: $\bm{y = x(y')^2 - \frac{1}{y'}}$} \\
$y' = p \implies y = xp^2 - \frac{1}{p}$.
\begin{align*}
dy &= p^2 dx + 2xp \, dp + \frac{1}{p^2}dp \qquad dy=pdx
\end{align*}
\begin{align*}
p dx &= p^2 dx + \left(2xp + \frac{1}{p^2}\right)dp \implies \frac{dx}{dp} - \frac{2}{1-p}x = \frac{1}{p^3(1-p)}
\end{align*}
\textit{Resolviendo la lineal, $x(p) = \dfrac{2Cp^2 + 2p - 1}{2p^2(1-p)^2}$. Además, $p=1 \implies y = x-1$ constituye una solución singular.}
\[
\boxed{\begin{cases} x = \dfrac{2Cp^2 + 2p - 1}{2p^2(1-p)^2} \\[2ex] y = \dfrac{2Cp^2 + 2p - 1}{2(1-p)^2} - \dfrac{1}{p} \end{cases}}
\]


\noindent\textbf{10. Integrad la EDO: $\bm{y = xy' + \frac{a}{(y')^2}}$} \\
$y' = p \implies y = xp + \frac{a}{p^2}$.
\begin{align*}
dy &= pdx + xdp - \frac{2a}{p^3}dp \implies pdx = pdx + \left(x - \frac{2a}{p^3}\right)dp \implies \left(x - \frac{2a}{p^3}\right)dp = 0
\end{align*}
\begin{itemize}
\item $dp = 0 \implies p = C \implies$ Solución general:
\[
\boxed{y = Cx + \frac{a}{C^2}}
\]
\item $x - \dfrac{2a}{p^3} = 0 \implies p = \left(\dfrac{2a}{x}\right)^{1/3} \implies y = 3\left(\dfrac{ax^2}{4}\right)^{1/3} \implies$ Solución singular:
\[
\boxed{4y^3 = 27ax^2}
\]
\end{itemize}


\noindent\textbf{11. Integrad la EDO: $\bm{x(y')^2 - yy' - y' + 1 = 0}$} \\
Hacemos $y' = p$:
\begin{align*}
xp^2 - yp - p + 1 = 0 \implies y = xp - 1 + \frac{1}{p}
\end{align*}
\textit{$dy = pdx + xdp - \frac{1}{p^2}dp \implies \left(x - \frac{1}{p^2}\right)dp = 0$.}
\begin{itemize}
\item $dp = 0 \implies p = C \implies$ Solución general:
\[
\boxed{y = Cx + \frac{1}{C} - 1}
\]
\item $x - \dfrac{1}{p^2} = 0 \implies p = \dfrac{1}{\sqrt{x}} \implies y = 2\sqrt{x} - 1 \implies$ Solución particular:
\[
\boxed{(y+1)^2 = 4x}
\]
\end{itemize}


\noindent\textbf{12. Integrad la EDO de Riccati: $\bm{y'e^{-x} + y^2 - 2ye^x = 1 - e^{2x}}$ sabiendo la solución $\bm{y_1 = e^x}$} \\
$y = y_1 + z = e^x + z \implies y' = e^x + z'$. Sustituyendo:
\begin{align*}
(e^x + z')e^{-x} + (e^x+z)^2 - 2(e^x+z)e^x &= 1 - e^{2x} \\
1 + z'e^{-x} + e^{2x} + 2ze^x + z^2 - 2e^{2x} - 2ze^x &= 1 - e^{2x} \implies z'e^{-x} + z^2 = 0
\end{align*}
\begin{align*}
\frac{dz}{z^2} &= -e^x dx \implies -\frac{1}{z} = -e^x - C \implies z = \frac{1}{e^x + C}
\end{align*}
\[
\boxed{y(x) = e^x + \frac{1}{e^x + C}}
\]


\noindent\textbf{13. Hallad la solución general de la EDO de Riccati cuando $\bm{a(x):b(x):c(x) = m:n:p}$} \\
Como los coeficientes son proporcionales a constantes $m,n,p$, podemos usar $\phi(x)$ tal que $a(x) = m\phi(x)$, $b(x) = n\phi(x)$ y $c(x) = p\phi(x)$. La EDO es de variables separables:
\begin{align*}
y' + \phi(x)\big(my^2 + ny + p\big) = 0 \implies \frac{dy}{my^2 + ny + p} = -\phi(x)dx
\end{align*}
Definimos el discriminante $\Delta = n^2 - 4mp$:
\begin{itemize}
\item \textbf{ $\bm{\Delta > 0}$:} Existen dos raíces reales distintas $r_1, r_2 = \frac{-n \pm \sqrt{\Delta}}{2m}$. \\
\begin{align*}
\frac{1}{m(y-r_1)(y-r_2)} &= \frac{A}{y-r_1} + \frac{B}{y-r_2} \implies 1 = m\big[A(y-r_2) + B(y-r_1)\big]
\end{align*}
\begin{align*}
A &= \frac{1}{m(r_1 - r_2)} = \frac{1}{\sqrt{\Delta}}, \qquad B = \frac{1}{m(r_2 - r_1)} = -\frac{1}{\sqrt{\Delta}}
\end{align*}
\begin{align*}
\int \left( \frac{1/\sqrt{\Delta}}{y-r_1} - \frac{1/\sqrt{\Delta}}{y-r_2} \right) dy &= -\int \phi(x)dx \\
\frac{1}{\sqrt{\Delta}}\left(\ln|y-r_1| - \ln|y-r_2|\right) &= -\int \phi(x)dx + C_1 \implies \frac{1}{\sqrt{\Delta}}\ln\left|\frac{y-r_1}{y-r_2}\right| = -\int \phi(x)dx + C_1
\end{align*}
\[
\boxed{\frac{y-r_1}{y-r_2} = K e^{-\sqrt{\Delta}\int \phi(x)dx}}
\]

\item \textbf{$\bm{\Delta = 0}$:} raíz real doble $r = -\frac{n}{2m}$.
\begin{align*}
\int \frac{dy}{m(y + \frac{n}{2m})^2} &= -\int \phi(x)dx \implies -\frac{2}{2my+n} = -\int \phi(x)dx - C
\end{align*}
\[
\boxed{\frac{2}{2my+n} = \int \phi(x)dx + C}
\]

\item \textbf{Caso $\bm{\Delta < 0}$:} No existen raíces reales. Completando el cuadrado en el denominador:
\begin{align*}
my^2 + ny + p = m\left[\left(y + \frac{n}{2m}\right)^2 + \frac{4mp-n^2}{4m^2}\right] = m\left[\left(y + \frac{n}{2m}\right)^2 + \frac{-\Delta}{4m^2}\right]
\end{align*}
Metiendo la integral en calculadora de integrales:
\begin{align*}
\int \frac{dy}{m\left[\left(y + \frac{n}{2m}\right)^2 + \left(\frac{\sqrt{-\Delta}}{2m}\right)^2\right]} &= -\int \phi(x)dx \\
\frac{2}{\sqrt{-\Delta}}\arctan\left(\frac{2my+n}{\sqrt{-\Delta}}\right) &= -\int \phi(x)dx + C_1
\end{align*}
\[
\boxed{\arctan\left(\frac{2my+n}{\sqrt{-\Delta}}\right) = \frac{\sqrt{-\Delta}}{2}\left(C - \int \phi(x)dx\right)}
\]
\end{itemize}

\noindent\textbf{14. Construid la EDO de la familia de curvas: $\bm{y^2 = 2Cx + C^2}$} \\
\begin{align*}
2yy' = 2C \implies C = yy'
\end{align*}
\begin{align*}
y^2 = 2x(yy') + (yy')^2 \implies y^2 = 2xyy' + y^2(y')^2
\end{align*}
Dividiendo entre $y$ ($y \neq 0$):
\[
\boxed{y(y')^2 + 2xy' - y = 0}
\]


\noindent\textbf{15. Construid la EDO de la familia de curvas: $\bm{y = ax^2 + bx + c}$} \\
Como la familia tiene 3 parámetros constantes ($a, b, c$), derivamos 3 veces:
\begin{align*}
y' &= 2ax + b \\
y'' &= 2a \\
y''' &= 0
\end{align*}
\[
\boxed{y''' = 0}
\]


\noindent\textbf{16. Construid la EDO de la familia de curvas: $\bm{y = C_1 x + C_2/x + C_3}$} \\
Como vimos en clase:
\begin{align*}
y' &= C_1 - \frac{C_2}{x^2} \\
y'' &= \frac{2C_2}{x^3} \implies C_2 = \frac{x^3 y''}{2} \\
y''' &= -\frac{6C_2}{x^4}
\end{align*}
\begin{align*}
y''' = -\frac{6}{x^4}\left(\frac{x^3 y''}{2}\right) = -\frac{3y''}{x} \implies xy''' + 3y'' = 0
\end{align*}
\[
\boxed{xy''' + 3y'' = 0}
\]


\noindent\textbf{17. Hallad las trayectorias ortogonales a la familia de curvas: $\bm{y^2 + 2ax = 0}$ con $\bm{a>0}$} \\
Despejamos $a = -\frac{y^2}{2x}$ y derivamos la familia $2yy' + 2a = 0 \implies a = -yy'$.
Igualando ambas expresiones de $a$ obtenemos la EDO de la familia original:
\begin{align*}
-yy' = -\frac{y^2}{2x} \implies y' = \frac{y}{2x}
\end{align*}
$y \to \tilde{y}$ y $y' \to -\frac{1}{\tilde{y}'}$:
\begin{align*}
-\frac{1}{\tilde{y}'} = \frac{\tilde{y}}{2x} \implies \tilde{y}' = -\frac{2x}{\tilde{y}} \implies \tilde{y} \, d\tilde{y} = -2x \, dx
\end{align*}
$\frac{\tilde{y}^2}{2} = -x^2 + C_1 \implies 2x^2 + \tilde{y}^2 = C$.
\[
\boxed{2x^2 + \tilde{y}^2 = C, \qquad C > 0}
\]


\noindent\textbf{18. Hallad las trayectorias ortogonales a la familia de curvas: $\bm{y = ax^n}$ con $\bm{n}$ fijo} \\
Despejamos $a = \frac{y}{x^n}$ y derivamos $y' = nax^{n-1} = n\left(\frac{y}{x^n}\right)x^{n-1} = \frac{ny}{x}$. \\
EDO de la familia original: $y' = \frac{ny}{x}$. \\
$y \to \tilde{y}$ y $y' \to -\frac{1}{\tilde{y}'}$:
\begin{align*}
-\frac{1}{\tilde{y}'} = \frac{n\tilde{y}}{x} \implies \tilde{y}' = -\frac{x}{n\tilde{y}} \implies n\tilde{y} \, d\tilde{y} = -x \, dx
\end{align*}
$\frac{n\tilde{y}^2}{2} = -\frac{x^2}{2} + C_1$.
\[
\boxed{x^2 + n\tilde{y}^2 = C}
\]


\noindent\textbf{19. Hallad las trayectorias ortogonales a la familia de curvas: $\bm{y = ae^{kx}}$ con $\bm{k}$ fijo} \\
$y' = kae^{kx} = ky$. La EDO de la familia es $y' = ky$. \\
$y \to \tilde{y}$ y $y' \to -\frac{1}{\tilde{y}'}$:
\begin{align*}
-\frac{1}{\tilde{y}'} = k\tilde{y} \implies \tilde{y}' = -\frac{1}{k\tilde{y}} \implies k\tilde{y} \, d\tilde{y} = -dx
\end{align*}
$\frac{k\tilde{y}^2}{2} = -x + C_1 \implies k\tilde{y}^2 + 2x = C$.
\[
\boxed{k\tilde{y}^2 + 2x = C}
\]


\noindent\textbf{20. Hallad las trayectorias ortogonales a la familia de curvas: $\bm{\cos y = ae^{-x}}$} \\
$-\sin y \cdot y' = -ae^{-x} = -\cos y \implies y' = \cot y$. \\
$y \to \tilde{y}$ y $y' \to -\frac{1}{\tilde{y}'}$:
\begin{align*}
-\frac{1}{\tilde{y}'} = \cot \tilde{y} \implies \tilde{y}' = -\tan \tilde{y} \implies \frac{d\tilde{y}}{\tan \tilde{y}} = -dx \implies \cot \tilde{y} \, d\tilde{y} = -dx
\end{align*}
Integrando: $\ln|\sin \tilde{y}| = -x + C_1 \implies |\sin \tilde{y}| = e^{-x+C_1} = C e^{-x}$.
\[
\boxed{\sin \tilde{y} = C e^{-x}}
\]
\separador
\section{Hoja 4}

\noindent\textbf{1. Verificar que $\bm{y_1=1}$ e $\bm{y_2=\ln x}$ son soluciones de la EDO $\bm{xy''+y'=0}$ y escribir la solución general.} \\
$y_1=1 \implies y_1'=0, \, y_1''=0$:
\begin{align*}
x(0) + 0 = 0 \quad \checkmark
\end{align*}
$y_2=\ln x \implies y_2'=\frac{1}{x}, \, y_2''=-\frac{1}{x^2}$:
\begin{align*}
x\left(-\frac{1}{x^2}\right) + \frac{1}{x} = -\frac{1}{x} + \frac{1}{x} = 0 \quad \checkmark
\end{align*}
$y_1$ e $y_2$ son li., luego
\[
\boxed{y(x) = C_1 + C_2 \ln x}
\]


\noindent\textbf{2. Hallad la EDO de la familia de curvas: $\bm{y=c_1 x + c_2 x^2}$} \\
Derivamos sucesivamente dos veces respecto de $x$:
\begin{align*}
y' &= c_1 + 2c_2 x \\
y'' &= 2c_2 \implies c_2 = \frac{y''}{2} \implies c_1 = y' - 2c_2 x = y' - x y''
\end{align*}
Despejamos
\begin{align*}
y = (y' - x y'')x + \left(\frac{y''}{2}\right)x^2 = x y' - \frac{1}{2}x^2 y'' \implies x^2 y'' - 2x y' + 2y = 0
\end{align*}
\[
\boxed{x^2 y'' - 2x y' + 2y = 0}
\]


\noindent\textbf{3. Hallad la EDO de la familia de curvas: $\bm{y=c_1 e^{kx} + c_2 e^{-kx}}$} \\
\begin{align*}
y' &= k c_1 e^{kx} - k c_2 e^{-kx} \\
y'' &= k^2 c_1 e^{kx} + k^2 c_2 e^{-kx} = k^2 (c_1 e^{kx} + c_2 e^{-kx}) = k^2 y
\end{align*}
\[
\boxed{y'' - k^2 y = 0}
\]


\noindent\textbf{4. Hallad la EDO de la familia de curvas: $\bm{y=c_1 x + c_2 \sin x}$} \\
\begin{align*}
y' &= c_1 + c_2 \cos x \\
y'' &= -c_2 \sin x \implies c_2 = -\frac{y''}{\sin x} \\
c_1 x &= y - c_2 \sin x = y + y'' \implies c_1 = \frac{y + y''}{x} \\
\end{align*}
\begin{align*}
y' = \frac{y + y''}{x} - y'' \cot x \implies x \sin x \cdot y' = \sin x (y + y'') - x \cos x \cdot y''
\end{align*}
Reordenando términos:
\[
\boxed{(x \cos x - \sin x) y'' + x \sin x \cdot y' - y \sin x = 0}
\]


\noindent\textbf{5. Si $\bm{y_1}$ e $\bm{y_2}$ son dos soluciones de una EDO lineal homogénea de segundo orden sobre $\bm{[a,b]}$ que tienen un cero común en ese intervalo, demostrad que una es múltiplo de la otra.} \\
Sea $x_0 \in [a,b]$ tal que $y_1(x_0) = y_2(x_0) = 0$. Evaluamos el Wronskiano en $x_0$:
\begin{align*}
W(y_1, y_2)(x_0) = y_1(x_0)y_2'(x_0) - y_1'(x_0)y_2(x_0) = 0 \cdot y_2'(x_0) - y_1'(x_0) \cdot 0 = 0
\end{align*}
Por el Lema de clase, si el Wronskiano de dos soluciones se anula en un punto de $[a,b]$, es idénticamente nulo en todo el intervalo: $W(y_1, y_2)(x) \equiv 0 \quad \forall x \in [a,b]$.\\
\textit{Un Wronskiano idénticamente nulo implica que $y_1$ e $y_2$ son linealmente dependientes en $[a,b]$, por lo que existe una constante $C$ tal que $y_2(x) = C y_1(x)$ o $y_1(x) = C y_2(x)$.} \quad $\blacksquare$\\


\noindent\textbf{6. Demostrad que $\bm{e^x}$ y $\bm{e^{-x}}$ son soluciones linealmente independientes de $\bm{y''-y=0}$ sobre cualquier intervalo cerrado.} \\
$y_1 = e^x \implies y_1'' = e^x \implies e^x - e^x = 0$, y $y_2 = e^{-x} \implies y_2'' = e^{-x} \implies e^{-x} - e^{-x} = 0$. Ambos satisfacen la EDO. \\
Calculamos el Wronskiano:
\begin{align*}
W(e^x, e^{-x}) = \begin{vmatrix} e^x & e^{-x} \\ e^x & -e^{-x} \end{vmatrix} = e^x(-e^{-x}) - e^x(e^{-x}) = -1 - 1 = -2 \neq 0 \quad \forall x \in \mathbb{R}
\end{align*}
\textit{Como $W \neq 0$ en todo intervalo cerrado $[a,b]$, $e^x$ y $e^{-x}$ son li.} \quad $\blacksquare$\\


\noindent\textbf{7. Comprobad que las funciones $\bm{y_1}$ e $\bm{y_2}$ son soluciones linealmente independientes de la EDO dada sobre $\bm{[0,2]}$ y hallad la solución particular:}

\begin{itemize}
\item \textbf{(a) $\bm{y''+y'-2y=0}$, $\bm{y_1=e^x}$, $\bm{y_2=e^{-2x}}$, $\bm{y(0)=8, y'(0)=2}$:}
\begin{align*}
W(e^x, e^{-2x}) = \begin{vmatrix} e^x & e^{-2x} \\ e^x & -2e^{-2x} \end{vmatrix} = -3e^{-x} \neq 0 \quad \forall x \in [0,2]
\end{align*}
$y(x) = C_1 e^x + C_2 e^{-2x}$.Con $y(0)=8 \implies C_1 + C_2 = 8$, y $y'(0)=2 \implies C_1 - 2C_2 = 2$.

\[
\boxed{y(x) = 6e^x + 2e^{-2x}}
\]

\item \textbf{(b) $\bm{y''+y'-2y=0}$, $\bm{y_1=e^x}$, $\bm{y_2=e^{-2x}}$, $\bm{y(1)=0, y'(1)=0}$:}
\begin{align*}
\begin{cases} C_1 e + C_2 e^{-2} = 0 \\ C_1 e - 2C_2 e^{-2} = 0 \end{cases} \implies 3C_2 e^{-2} = 0 \implies C_2 = 0, C_1 = 0
\end{align*}
\[
\boxed{y(x) = 0}
\]

\item \textbf{(c) $\bm{y''+5y'+6y=0}$, $\bm{y_1=e^{-2x}}$, $\bm{y_2=e^{-3x}}$, $\bm{y(0)=1, y'(0)=1}$:}\\
$y(x) = C_1 e^{-2x} + C_2 e^{-3x}$.
\begin{align*}
\begin{cases} C_1 + C_2 = 1 \\ -2C_1 - 3C_2 = 1 \end{cases} \implies C_1 = 4, \quad C_2 = -3
\end{align*}
\[
\boxed{y(x) = 4e^{-2x} - 3e^{-3x}}
\]

\item \textbf{(d) $\bm{y''+y'=0}$, $\bm{y_1=1}$, $\bm{y_2=e^{-x}}$, $\bm{y(2)=0, y'(2)=1/e^2}$:}
\begin{align*}
W(1, e^{-x}) = -e^{-x} \neq 0 \quad \forall x \in [0,2]
\end{align*}
$y(x) = C_1 + C_2 e^{-x} \implies y'(x) = -C_2 e^{-x}$.\\
$y'(2) = -C_2 e^{-2} = e^{-2} \implies C_2 = -1$. \\
$y(2) = C_1 - e^{-2} = 0 \implies C_1 = e^{-2}$.\\
\[
\boxed{y(x) = e^{-2} - e^{-x}}
\]
\end{itemize}


\noindent\textbf{8. Verificad que $\bm{y_1(x)=\sin x}$ es solución de $\bm{y''+y=0}$. Obtened otra solución $\bm{y_2(x)}$ L.I. y la solución general.} \\
Comprobamos la primera solución derivando: $(\sin x)'' + \sin x = -\sin x + \sin x = 0 \quad \checkmark$. \\
La EDO es lineal homogénea con coef const,
\begin{align*}
s^2 + 1 = 0 \implies s = \pm i
\end{align*}
Las raíces imaginarias puras generan soluciones trigonométricas. Como ya conocemos $y_1(x) = \sin x$, la otra solución linealmente independiente es:
\[
y_2(x) = \cos x
\]
Por lo tanto, la solución general es:
\[
\boxed{y(x) = C_1 \sin x + C_2 \cos x}
\]


\noindent\textbf{9. La EDO $\bm{xy''+3y'=0}$ admite $\bm{y_1(x)=1}$. Hallad otra solución $\bm{y_2(x)}$ L.I. sobre cualquier intervalo que no contenga $\bm{0}$ y la solución general.} \\
$p = y' \implies p' = y''$:
\begin{align*}
x p' + 3p = 0 \implies x \frac{dp}{dx} = -3p \implies \frac{dp}{p} = -\frac{3}{x} dx
\end{align*}
\begin{align*}
\ln|p| = -3\ln|x| + C \implies p(x) = \frac{K}{x^3}
\end{align*}
\begin{align*}
y(x) = \int \frac{K}{x^3} dx = -\frac{K}{2x^2} + C_1 = \frac{C_2}{x^2} + C_1
\end{align*}
\begin{align*}
W(1, x^{-2}) = -2x^{-3} \neq 0 \quad \text{para } x \neq 0
\end{align*}
\[
\boxed{y(x) = C_1 + \frac{C_2}{x^2}}
\]


\noindent\textbf{10. Comprobad que $\bm{y_1(x)=e^x}$ es solución de $\bm{xy''-(2x+1)y'+(x+1)y=0}$ y hallad la solución general.} \\
Sustituyendo $y_1 = e^x \implies y_1' = e^x, \, y_1'' = e^x$:
\begin{align*}
x e^x - (2x+1)e^x + (x+1)e^x = e^x(x - 2x - 1 + x + 1) = 0 \quad \checkmark
\end{align*}
$y_2(x) = v(x) e^x \implies y_2' = (v'+v)e^x, \, y_2'' = (v''+2v'+v)e^x$.
\begin{align*}
x(v''+2v'+v)e^x - (2x+1)(v'+v)e^x + (x+1)v e^x &= 0 \\
x v'' + [2x - (2x+1)] v' &= 0 \implies x v'' - v' = 0 \implies \frac{v''}{v'} = \frac{1}{x}
\end{align*}
$\ln v' = \ln x \implies v' = x \implies v(x) = \frac{x^2}{2} \implies y_2(x) = x^2 e^x$.
\[
\boxed{y(x) = C_1 e^x + C_2 x^2 e^x}
\]


\noindent\textbf{11. Hallad la solución general de las siguientes ecuaciones:}

\begin{itemize}
\item \textbf{(a) $\bm{y''-4y'+4y=0}$:} $(s-2)^2 = 0 \implies s = 2$ (doble). $\quad \boxed{y(x) = (C_1 + C_2 x)e^{2x}}$
\item \textbf{(b) $\bm{y''-9y'+20y=0}$:} $(s-4)(s-5) = 0 \implies s_1 = 4, s_2 = 5$. $\quad \boxed{y(x) = C_1 e^{4x} + C_2 e^{5x}}$
\item \textbf{(c) $\bm{2y''+2y'+3y=0}$:} $s = -\frac{1}{2} \pm i\frac{\sqrt{5}}{2}$. $\quad \boxed{y(x) = e^{-x/2}\left(C_1 \cos\frac{\sqrt{5}}{2}x + C_2 \sin\frac{\sqrt{5}}{2}x\right)}$
\item \textbf{(d) $\bm{4y''-12y'+9y=0}$:} $(2s-3)^2 = 0 \implies s = 3/2$ (doble). $\quad \boxed{y(x) = (C_1 + C_2 x)e^{3x/2}}$
\item \textbf{(e) $\bm{y''+y'=0}$:} $s(s+1) = 0 \implies s_1 = 0, s_2 = -1$. $\quad \boxed{y(x) = C_1 + C_2 e^{-x}}$
\item \textbf{(f) $\bm{y''-6y'+25y=0}$:} $s = 3 \pm 4i$. $\quad \boxed{y(x) = e^{3x}(C_1 \cos 4x + C_2 \sin 4x)}$
\item \textbf{(g) $\bm{4y''+20y'+25y=0}$:} $(2s+5)^2 = 0 \implies s = -5/2$ (doble). $\quad \boxed{y(x) = (C_1 + C_2 x)e^{-5x/2}}$
\item \textbf{(h) $\bm{y''+2y'+3y=0}$:} $s = -1 \pm i\sqrt{2}$. $\quad \boxed{y(x) = e^{-x}(C_1 \cos\sqrt{2}x + C_2 \sin\sqrt{2}x)}$
\item \textbf{(i) $\bm{y''=4y}$:} $s^2 - 4 = 0 \implies s = \pm 2$. $\quad \boxed{y(x) = C_1 e^{2x} + C_2 e^{-2x}}$
\item \textbf{(j) $\bm{4y''-8y'+7y=0}$:} $s = 1 \pm i\frac{\sqrt{3}}{2}$. $\quad \boxed{y(x) = e^x\left(C_1 \cos\frac{\sqrt{3}}{2}x + C_2 \sin\frac{\sqrt{3}}{2}x\right)}$
\item \textbf{(k) $\bm{2y''+y'-y=0}$:} $(2s-1)(s+1) = 0 \implies s_1 = 1/2, s_2 = -1$. $\quad \boxed{y(x) = C_1 e^{x/2} + C_2 e^{-x}}$
\item \textbf{(l) $\bm{16y''-8y'+y=0}$:} $(4s-1)^2 = 0 \implies s = 1/4$ (doble). $\quad \boxed{y(x) = (C_1 + C_2 x)e^{x/4}}$
\item \textbf{(m) $\bm{y''+4y'+5y=0}$:} $s = -2 \pm i$. $\quad \boxed{y(x) = e^{-2x}(C_1 \cos x + C_2 \sin x)}$
\item \textbf{(n) $\bm{y''+4y'-5y=0}$:} $(s+5)(s-1) = 0 \implies s_1 = 1, s_2 = -5$. $\quad \boxed{y(x) = C_1 e^x + C_2 e^{-5x}}$\\
\end{itemize}

\noindent\textbf{12. Resolved los siguientes problemas con condiciones iniciales:}

\begin{itemize}
\item \textbf{(a) $\bm{y''-5y'+6y=0}$ con $\bm{y(1)=e^2, y'(1)=3e^2}$:} \\
$y(x) = C_1 e^{2x} + C_2 e^{3x}$. Evaluando en $x=1$:
\begin{align*}
\begin{cases} C_1 e^2 + C_2 e^3 = e^2 \\ 2C_1 e^2 + 3C_2 e^3 = 3e^2 \end{cases} \implies C_1 = 0, \quad C_2 = e^{-1}
\end{align*}
\[
\boxed{y(x) = e^{3x-1}}
\]

\item \textbf{(b) $\bm{y''-6y'+5y=0}$ con $\bm{y(0)=3, y'(0)=11}$:} \\
$y(x) = C_1 e^x + C_2 e^{5x}$.
\begin{align*}
\begin{cases} C_1 + C_2 = 3 \\ C_1 + 5C_2 = 11 \end{cases} \implies 4C_2 = 8 \implies C_2 = 2, \quad C_1 = 1
\end{align*}
\[
\boxed{y(x) = e^x + 2e^{5x}}
\]

\item \textbf{(c) $\bm{y''-6y'+9y=0}$ con $\bm{y(0)=0, y'(0)=5}$:} \\
$y(x) = (C_1 + C_2 x)e^{3x} \implies y(0) = C_1 = 0$. $y'(x) = C_2 e^{3x} + 3C_2 x e^{3x} \implies y'(0) = C_2 = 5$.
\[
\boxed{y(x) = 5x e^{3x}}
\]

\item \textbf{(d) $\bm{y''+4y'+5y=0}$ con $\bm{y(0)=1, y'(0)=0}$:} \\
$y(x) = e^{-2x}(C_1 \cos x + C_2 \sin x) \implies y(0) = C_1 = 1$. \\
$y'(0) = -2C_1 + C_2 = 0 \implies C_2 = 2$.
\[
\boxed{y(x) = e^{-2x}(\cos x + 2\sin x)}
\]

\item \textbf{(e) $\bm{y''+4y'+2y=0}$ con $\bm{y(0)=-1, y'(0)=2+3\sqrt{2}}$:} \\
$s = -2 \pm \sqrt{2} \implies y(x) = C_1 e^{(-2+\sqrt{2})x} + C_2 e^{(-2-\sqrt{2})x}$.
\begin{align*}
\begin{cases} C_1 + C_2 = -1 \\ (-2+\sqrt{2})C_1 + (-2-\sqrt{2})C_2 = 2+3\sqrt{2} \end{cases} \implies C_1 = 1, \quad C_2 = -2
\end{align*}
\[
\boxed{y(x) = e^{(-2+\sqrt{2})x} - 2e^{(-2-\sqrt{2})x}}
\]

\item \textbf{(f) $\bm{y''+8y'-9y=0}$ con $\bm{y(1)=2, y'(1)=0}$:} \\
$y(x) = C_1 e^x + C_2 e^{-9x}$.
\begin{align*}
\begin{cases} C_1 e + C_2 e^{-9} = 2 \\ C_1 e - 9C_2 e^{-9} = 0 \end{cases} \implies 10C_2 e^{-9} = 2 \implies C_2 = \frac{e^9}{5}, \quad C_1 = \frac{9}{5e}
\end{align*}
\[
\boxed{y(x) = \frac{9}{5}e^{x-1} + \frac{1}{5}e^{-9(x-1)}}
\]
\end{itemize}

\noindent\textbf{13. Ecuación equidimensional de Euler: $\bm{x^2y''+pxy'+qy=0}$. Demostrad que $\bm{x=e^z}$ la transforma en coeficientes constantes y aplicadla:} \\
Demostración: $z = \ln x \implies \frac{dz}{dx} = \frac{1}{x}$.
\begin{align*}
y' &= \frac{dy}{dz}\frac{dz}{dx} = \frac{1}{x}\frac{dy}{dz} \implies x y' = \frac{dy}{dz} \\
y'' &= \frac{d}{dx}\left(\frac{1}{x}\frac{dy}{dz}\right) = -\frac{1}{x^2}\frac{dy}{dz} + \frac{1}{x^2}\frac{d^2y}{dz^2} \implies x^2 y'' = \frac{d^2y}{dz^2} - \frac{dy}{dz}
\end{align*}
Sustituyendo en la EDO: $\frac{d^2y}{dz^2} + (p-1)\frac{dy}{dz} + qy = 0$, con coeficientes constantes. \quad $\blacksquare$

\begin{itemize}
\item \textbf{(a) $\bm{4x^2y''-3y=0}$:} $s(s-1) - 3/4 = 0 \implies s_1 = 3/2, s_2 = -1/2$. $\quad \boxed{y(x) = C_1 x^{3/2} + C_2 x^{-1/2}}$
\item \textbf{(b) $\bm{x^2y''-3xy'+4y=0}$:} $s(s-1) - 3s + 4 = (s-2)^2 = 0$. $\quad \boxed{y(x) = (C_1 + C_2 \ln x)x^2}$
\item \textbf{(c) $\bm{x^2y''+2xy'-6y=0}$:} $s(s-1) + 2s - 6 = (s+3)(s-2) = 0$. $\quad \boxed{y(x) = C_1 x^2 + C_2 x^{-3}}$
\item \textbf{(d) $\bm{x^2y''+2xy'+3y=0}$:} $s^2 + s + 3 = 0 \implies s = -\frac{1}{2} \pm i\frac{\sqrt{11}}{2}$. \\
$\boxed{y(x) = x^{-1/2}\left(C_1 \cos\left(\frac{\sqrt{11}}{2}\ln x\right) + C_2 \sin\left(\frac{\sqrt{11}}{2}\ln x\right)\right)}$
\item \textbf{(e) $\bm{x^2y''+xy'-2y=0}$:} $s^2 - 2 = 0 \implies s = \pm\sqrt{2}$. $\quad \boxed{y(x) = C_1 x^{\sqrt{2}} + C_2 x^{-\sqrt{2}}}$
\item \textbf{(f) $\bm{x^2y''+xy'-16y=0}$:} $s^2 - 16 = 0 \implies s = \pm 4$. $\quad \boxed{y(x) = C_1 x^4 + C_2 x^{-4}}$
\end{itemize}

\noindent\textbf{14. Demostrad que $\bm{y''+P(x)y'+Q(x)y=0}$ se transforma en coeficientes constantes con $\bm{z(x)}$ s.o.s. $\bm{(Q'+2PQ)/Q^{3/2}}$ es constante.} \\
Haciendo la transformación $z = z(x)$:
\begin{align*}
\frac{dy}{dx} = z' \frac{dy}{dz}, \qquad \frac{d^2y}{dx^2} = (z')^2 \frac{d^2y}{dz^2} + z'' \frac{dy}{dz}
\end{align*}
Sustituyendo en la EDO y dividiendo entre $(z')^2$:
\begin{align*}
\frac{d^2y}{dz^2} + \frac{z'' + P(x)z'}{(z')^2} \frac{dy}{dz} + \frac{Q(x)}{(z')^2} y = 0
\end{align*}
Denominamos a los nuevos coef (constantes) $p = \dfrac{z'' + P(x)z'}{(z')^2}$ y $q = \dfrac{Q(x)}{(z')^2}$, de modo que la ecuación adopta la forma:
\begin{align*}
\frac{d^2y}{dz^2} + p\frac{dy}{dz} + qy = 0
\end{align*}
\begin{align*}
\frac{Q(x)}{(z')^2} = q \implies z'(x) = \sqrt{q Q(x)} \implies z = \int \sqrt{qQ(x)} dx
\end{align*}
Podemos elegir sin pérdida $q= 1$
\begin{align*}
Q(x) = (z')^2 \implies z'(x) = \sqrt{Q(x)} \implies z = \int \sqrt{Q(x)} dx
\end{align*}
Derivando $z'(x)$ obtenemos $z'' = \frac{Q'(x)}{2\sqrt{Q(x)}}$. Sustituyendo esto en el coeficiente $p$:
\begin{align*}
p = \frac{\frac{Q'(x)}{2\sqrt{Q(x)}} + P(x)\sqrt{Q(x)}}{Q(x)} = \frac{Q'(x) + 2P(x)Q(x)}{2[Q(x)]^{3/2}}
\end{align*}
\begin{align*}
\frac{Q'(x) + 2P(x)Q(x)}{[Q(x)]^{3/2}} = C \quad (\text{constante})
\end{align*}
Quedando finalmente reducida a $\frac{d^2y}{dz^2} + C \frac{dy}{dz} + y = 0$. \qquad $\blacksquare$ \\

\noindent\textbf{15. Usad el resultado del problema anterior y resolved cuando sea posible:}

\begin{itemize}
\item \textbf{(a) $\bm{xy''+(x^2-1)y'+x^3y=0} \implies y'' + \left(x - \frac{1}{x}\right)y' + x^2 y = 0$:} \\
$P(x) = x - 1/x, \, Q(x) = x^2 \implies \frac{Q'+2PQ}{Q^{3/2}} = \frac{2x + 2(x-1/x)x^2}{x^3} = 2$ (constante). \\
$z = \int x \, dx = \frac{x^2}{2} \implies \frac{d^2y}{dz^2} + \frac{dy}{dz} + y = 0 \implies s = -\frac{1}{2} \pm i\frac{\sqrt{3}}{2}$.
\[
\boxed{y(x) = e^{-x^2/4}\left(C_1 \cos\frac{\sqrt{3}x^2}{4} + C_2 \sin\frac{\sqrt{3}x^2}{4}\right)}
\]

\item \textbf{(b) $\bm{y''+3xy'+x^2y=0}$:} \\
$P(x) = 3x, \, Q(x) = x^2 \implies \frac{Q'+2PQ}{2Q^{3/2}} = \frac{2x + 6x^3}{2x^3} = 3 + \frac{1}{x^2}$ (no es constante).
\[
\boxed{\text{No es posible transformarla en coeficientes constantes.}}
\]
\end{itemize}


\noindent\textbf{16. Hallad la solución general de las siguientes ecuaciones no homogéneas:}

\begin{itemize}
\item \textbf{(a) $\bm{y''-y'-6y=20e^{-2x}}$:} \\
$y_h = C_1 e^{3x} + C_2 e^{-2x}$. Proponemos $y_p = A x e^{-2x} \implies y_p'' - y_p' - 6y_p = -5A e^{-2x} = 20e^{-2x} \implies A = -4$.
\[
\boxed{y(x) = C_1 e^{3x} + C_2 e^{-2x} - 4x e^{-2x}}
\]

\item \textbf{(b) $\bm{y''-3y'+2y=14\sin 2x - 18\cos 2x}$:} \\
$y_h = C_1 e^x + C_2 e^{2x}$. Proponemos $y_p = A\cos 2x + B\sin 2x$. Sustituyendo e igualando:
\begin{align*}
(-2A - 6B)\cos 2x + (6A - 2B)\sin 2x = -18\cos 2x + 14\sin 2x \implies A = 3, \quad B = 2
\end{align*}
\[
\boxed{y(x) = C_1 e^x + C_2 e^{2x} + 3\cos 2x + 2\sin 2x}
\]

\item \textbf{(c) $\bm{y''+y=2\cos x}$:} \\
$y_h = C_1 \cos x + C_2 \sin x$. Por doble cero, $y_p = x(A\cos x + B\sin x) \implies y_p'' + y_p = -2A\sin x + 2B\cos x = 2\cos x \implies A = 0, B = 1$.
\[
\boxed{y(x) = C_1 \cos x + C_2 \sin x + x \sin x}
\]

\item \textbf{(d) $\bm{y''-2y'=12x-10}$:} \\
$y_h = C_1 + C_2 e^{2x}$. Proponemos $y_p = Ax^2 + Bx \implies y_p'' - 2y_p' = -4Ax + (2A - 2B) = 12x - 10 \implies A = -3, B = 2$.
\[
\boxed{y(x) = C_1 + C_2 e^{2x} - 3x^2 + 2x}
\]

\item \textbf{(e) $\bm{y''-2y'+y=6e^x}$:} \\
$y_h = (C_1 + C_2 x)e^x$. $y_p = A x^2 e^x \implies y_p'' - 2y_p' + y_p = 2A e^x = 6e^x \implies A = 3$.
\[
\boxed{y(x) = (C_1 + C_2 x + 3x^2)e^x}
\]

\item \textbf{(f) $\bm{y''-2y'+2y=e^x \sin x}$:} \\
$y_h = e^x(C_1 \cos x + C_2 \sin x)$. Por doble cero, $y_p = x e^x (A\cos x + B\sin x) \implies A = -1/2, B = 0$.
\[
\boxed{y(x) = e^x(C_1 \cos x + C_2 \sin x) - \frac{1}{2}x e^x \cos x}
\]

\item \textbf{(g) $\bm{y''+y'=10x^4+2}$:} \\
$y_h = C_1 + C_2 e^{-x}$. Proponemos (porque la edo no tiene término con y) $y_p = x(Ax^4 + Bx^3 + Cx^2 + Dx + E)   \implies A=2, B=-10, C=40, D=-120, E=242$.
\[
\boxed{y(x) = C_1 + C_2 e^{-x} + 2x^5 - 10x^4 + 40x^3 - 120x^2 + 242x}
\]

\item \textbf{(h) $\bm{y''+9y=2\sin 3x + 4\sin x - 26e^{-2x} + 27x^3}$:} \\
$s^2 + 9 = 0 \implies s = \pm 3i$,
\begin{align*}
y_h(x) = C_1 \cos 3x + C_2 \sin 3x
\end{align*}
Dividimos el término no homogéneo en cuatro partes:
\begin{itemize}
\item \textbf{Para $\bm{f_1(x) = 2\sin 3x}$ (Cero de la caracteristica):} proponemos $y_{p1} = x(A\cos 3x + B\sin 3x)$. Sustituyendo: $-6A\sin 3x = 2\sin 3x \implies A = -1/3, B = 0$.
\[
y_{p1}(x) = -\frac{1}{3}x\cos 3x
\]
\item \textbf{Para $\bm{f_2(x) = 4\sin x}$:} proponemos $y_{p2} = C\cos x + D\sin x$. Sustituyendo se llega a $8D\sin x = 4\sin x \implies C = 0, D = 1/2$.
\[
y_{p2}(x) = \frac{1}{2}\sin x
\]
\item \textbf{Para $\bm{f_3(x) = -26e^{-2x}}$:} proponemos $y_{p3} = Ee^{-2x}$. Sustituyendo se obtiene $13Ee^{-2x} = -26e^{-2x} \implies E = -2$.
\[
y_{p3}(x) = -2e^{-2x}
\]
\item \textbf{Para $\bm{f_4(x) = 27x^3}$:} proponemos un polinomio general de tercer grado $y_{p4} = Fx^3 + Gx^2 + Hx + I$. Al sustituir e igualar coeficientes se determina que $F=3, G=0, H=-2, I=0$.
\[
y_{p4}(x) = 3x^3 - 2x
\]
\end{itemize}
Sumando la solución homogénea y todas las soluciones particulares obtenidas:
\[
\boxed{y(x) = C_1 \cos 3x + C_2 \sin 3x - \frac{1}{3}x\cos 3x + \frac{1}{2}\sin x - 2e^{-2x} + 3x^3 - 2x}
\]
\end{itemize}


\noindent\textbf{17. Hallad una solución particular por variación de parámetros:} \\
Antes de resolver, se deduce el método visto en clase para la EDO general $y'' + P(x)y' + Q(x)y = f(x)$. variando las constantes:
\begin{align*}
y_p(x) = u_1(x)y_1(x) + u_2(x)y_2(x)
\end{align*}
Derivando e imponiendo la condición $u_1'y_1 + u_2'y_2 = 0$ para evitar derivadas de orden superior en las incógnitas, obtenemos el sistema:
\begin{align*}
\begin{cases} u_1'y_1 + u_2'y_2 = 0 \\ u_1'y_1' + u_2'y_2' = f(x) \end{cases}
\end{align*}
\begin{align*}
\begin{pmatrix} y_1 & y_2 \\ y_1' & y_2' \end{pmatrix} \begin{pmatrix} u_1' \\ u_2' \end{pmatrix} = \begin{pmatrix} 0 \\ f(x) \end{pmatrix}
\end{align*}
Invertimos la matriz de coeficientes (det es el Wronskiano $W = y_1y_2' - y_1'y_2 \neq 0$):
\begin{align*}
\begin{pmatrix} u_1' \\ u_2' \end{pmatrix} = \frac{1}{W} \begin{pmatrix} y_2' & -y_2 \\ -y_1' & y_1 \end{pmatrix} \begin{pmatrix} 0 \\ f(x) \end{pmatrix} = \frac{1}{W} \begin{pmatrix} -y_2 f(x) \\ y_1 f(x) \end{pmatrix}
\end{align*}

\begin{align*}
u_1' = -\frac{y_2 f(x)}{W} \implies u_1(x) = -\int \frac{y_2 f(x)}{W} dx, \qquad u_2' = \frac{y_1 f(x)}{W} \implies u_2(x) = \int \frac{y_1 f(x)}{W} dx
\end{align*}

\begin{itemize}
\item \textbf{(a) $\bm{y''+4y=\tan 2x}$:} \\
$y''+4y=0 \implies s^2+4=0 \implies s=\pm 2i$. 
\begin{align*}
y_1(x) &= \cos(2x) \\
y_2(x) &= \sin(2x)
\end{align*}
\begin{align*}
W = \begin{vmatrix} \cos(2x) & \sin(2x) \\ -2\sin(2x) & 2\cos(2x) \end{vmatrix} = 2\cos^2(2x) + 2\sin^2(2x) = 2
\end{align*}
Aplicamos las fórmulas de variación de parámetros con $f(x) = \tan(2x)$:
\begin{align*}
u_1'(x) &= -\frac{y_2(x)f(x)}{W} = -\frac{\sin(2x)\tan(2x)}{2} = -\frac{1}{2}\frac{\sin^2(2x)}{\cos(2x)} = -\frac{1}{2}\frac{1-\cos^2(2x)}{\cos(2x)} = -\frac{1}{2}\big(\sec(2x) - \cos(2x)\big) \\
u_1(x) &= -\frac{1}{2}\int \big(\sec(2x) - \cos(2x)\big)dx = -\frac{1}{4}\ln|\sec(2x) + \tan(2x)| + \frac{1}{4}\sin(2x)
\end{align*}
\begin{align*}
u_2'(x) &= \frac{y_1(x)f(x)}{W} = \frac{\cos(2x)\tan(2x)}{2} = \frac{\sin(2x)}{2} \\
u_2(x) &= \frac{1}{2}\int \sin(2x)dx = -\frac{1}{4}\cos(2x)
\end{align*}
Sustituyendo en la solución particular $y_p(x) = u_1(x)y_1(x) + u_2(x)y_2(x)$:
\begin{align*}
y_p(x) &= \left( -\frac{1}{4}\ln|\sec(2x) + \tan(2x)| + \frac{1}{4}\sin(2x) \right)\cos(2x) + \left( -\frac{1}{4}\cos(2x) \right)\sin(2x) \\
y_p(x) &= -\frac{1}{4}\cos(2x)\ln|\sec(2x) + \tan(2x)| + \frac{1}{4}\sin(2x)\cos(2x) - \frac{1}{4}\sin(2x)\cos(2x)
\end{align*}
\[
\boxed{y_p(x) = -\frac{1}{4}\cos(2x)\ln|\sec(2x) + \tan(2x)|}
\]
\item \textbf{(b) $\bm{y''+2y'+y=e^{-x}\ln x}$:} \\
$y_1 = e^{-x}, y_2 = x e^{-x}, W = e^{-2x}$. $u_1' = -x\ln x \implies u_1 = -\frac{x^2}{2}\ln x + \frac{x^2}{4}$. $u_2' = \ln x \implies u_2 = x\ln x - x$.
\[
\boxed{y_p(x) = \frac{1}{2}x^2 e^{-x}\ln x - \frac{3}{4}x^2 e^{-x}}
\]
\item \textbf{(c) $\bm{y''-2y'-3y=64xe^{-x}}$:} \\
$y''-2y'-3y=0 \implies s^2-2s-3=0 \implies s_1=3, s_2=-1$. \\
Como el término no homogéneo es $64xe^{-x}$ y $s=-1$ es raíz, tenemos un caso similar a cuando es doble. Proponemos la solución particular multiplicada por $x$:
\begin{align*}
y_p(x) &= x(Ax+B)e^{-x} = (Ax^2+Bx)e^{-x} \\
y_p'(x) &= (-Ax^2 + (2A-B)x + B)e^{-x} \\
y_p''(x) &= (Ax^2 + (-4A+B)x + 2A-2B)e^{-x}
\end{align*}
Sustituimos en la EDO original y agrupamos por potencias de $x$:
\begin{align*}
y_p'' - 2y_p' - 3y_p &= e^{-x}\big[ (A+2A-3A)x^2 + (-4A+B-4A+2B-3B)x + (2A-2B-2B) \big] \\
&= e^{-x}\big[ 0x^2 - 8Ax + (2A-4B) \big] = 64xe^{-x}
\end{align*}
Igualando coeficientes obtenemos el sistema:
\begin{align*}
\begin{cases} -8A = 64 \implies A = -8 \\ 2A - 4B = 0 \implies 2(-8) - 4B = 0 \implies B = -4 \end{cases}
\end{align*}
\[
\boxed{y_p(x) = (-8x^2 - 4x)e^{-x}}
\]

\item \textbf{(d) $\bm{y''+2y'+5y=e^{-x}\sec 2x}$:} \\
Proponemos el cambio de variable dependiente $y(x) = z(x)e^{-x}$. Basicamente proque eso nos permite eliminar $e^{-x}$ que se convervara en las derivadas. Derivamos usando la regla del producto:
\begin{align*}
y' &= (z' - z)e^{-x} \\
y'' &= (z'' - 2z' + z)e^{-x}
\end{align*}
Sustituyendo en la EDO original y factorizando la exponencial:
\begin{align*}
(z'' - 2z' + z)e^{-x} + 2(z' - z)e^{-x} + 5z e^{-x} &= e^{-x}\sec 2x \\
e^{-x}(z'' - 2z' + z + 2z' - 2z + 5z) &= e^{-x}\sec 2x \implies z'' + 4z = \sec 2x
\end{align*}
Resolvemos esta nueva EDO por el método de variación. La homogénea es $z''+4z=0 \implies z_1 = \cos 2x, \, z_2 = \sin 2x$, con Wronskiano $W = 2$.
\begin{align*}
u_1' &= -\frac{\sin(2x)\sec(2x)}{2} = -\frac{1}{2}\tan(2x) \implies u_1 = -\frac{1}{2}\int \frac{\sin 2x}{\cos 2x}dx = \frac{1}{4}\ln|\cos 2x| \\
u_2' &= \frac{\cos(2x)\sec(2x)}{2} = \frac{1}{2} \implies u_2 = \int \frac{1}{2}dx = \frac{x}{2}
\end{align*}
La solución particular para la variable transformada $z$ es: 
\begin{align*}
z_p(x) = u_1 z_1 + u_2 z_2 = \frac{1}{4}\cos(2x)\ln|\cos 2x| + \frac{x}{2}\sin(2x)
\end{align*}
Deshaciendo el cambio de variable ($y_p = z_p e^{-x}$):
\[
\boxed{y_p(x) = e^{-x}\left(\frac{1}{4}\cos(2x)\ln|\cos 2x| + \frac{x}{2}\sin(2x)\right)}
\]
\item \textbf{(e) $\bm{2y''+3y'+y=e^{-3x}}$:}
\[
\boxed{y_p(x) = \frac{1}{10}e^{-3x}}
\]

\item \textbf{(f) $\bm{y''-3y'+2y=\frac{1}{1+e^{-x}}}$:} \\
$y_1 = e^x, y_2 = e^{2x}, W = e^{3x}$. Integrando las expresiones de $u_1'$ y $u_2'$:
\[
\boxed{y_p(x) = (e^x + e^{2x})\ln(1+e^{-x}) - e^x}
\]
\end{itemize}

\noindent\textbf{19. Demostrad que el método de variación de los parámetros aplicado a $\bm{y''+y=f(x)}$ conduce a $\bm{y_p(x)=\int_0^x f(t)\sin(x-t)dt}$. Hallad la fórmula para $\bm{y''+k^2y=f(x)}$ con $\bm{k>0}$.}

\noindent Para $y''+y=f(x)$, las soluciones homogéneas son $y_1(x) = \cos x$ e $y_2(x) = \sin x$ con $W=1$.
\begin{align*}
u_1(x) = -\int_0^x f(t)\sin t \, dt, \qquad u_2(x) = \int_0^x f(t)\cos t \, dt
\end{align*}
\begin{align*}
y_p(x) &= u_1(x)\cos x + u_2(x)\sin x = \int_0^x f(t) \big[\sin x \cos t - \cos x \sin t\big] dt = \int_0^x f(t)\sin(x-t) dt \quad \blacksquare
\end{align*}
Para $y''+k^2 y = f(x)$, $y_1(x) = \cos(kx)$ e $y_2(x) = \sin(kx)$ con $W = k$:
\begin{align*}
u_1(x) = -\frac{1}{k}\int_0^x f(t)\sin(kt) \, dt, \qquad u_2(x) = \frac{1}{k}\int_0^x f(t)\cos(kt) \, dt
\end{align*}
\[
\boxed{y_p(x) = \frac{1}{k}\int_0^x f(t)\sin(k(x-t)) \, dt}
\]

\noindent\textbf{18. Hallad una solución particular de $\bm{y''+y=f(x)}$ ($\bm{y_1=\cos x, y_2=\sin x, W=1}$):}

\begin{itemize}
\item \textbf{(a) $\bm{f(x)=\sec x}$:} $\quad \boxed{y_p(x) = \cos x \ln|\cos x| + x \sin x}$
\item \textbf{(b) $\bm{f(x)=\cot^2 x}$:} $\quad \boxed{y_p(x) = \cos x \ln|\csc x + \cot x| - 2}$
\item \textbf{(c) $\bm{f(x)=\cot 2x}$:} $\quad \boxed{y_p(x) = \frac{1}{2}\cos x \ln|\sec x + \tan x| - \frac{1}{2}\sin x \ln|\csc x + \cot x|}$
\item \textbf{(d) $\bm{f(x)=x\cos x}$:} $\quad \boxed{y_p(x) = \frac{x^2}{4}\sin x + \frac{x}{4}\cos x}$
\item \textbf{(e) $\bm{f(x)=\tan x}$:} $\quad \boxed{y_p(x) = -\cos x \ln|\sec x + \tan x|}$
\item \textbf{(f) $\bm{f(x)=\sec x\tan x}$:} $\quad \boxed{y_p(x) = x\cos x + \sin x \ln|\sec x|}$
\item \textbf{(g) $\bm{f(x)=\sec x\csc x}$:} $\quad \boxed{y_p(x) = -\cos x \ln|\sec x + \tan x| - \sin x \ln|\csc x + \cot x|}$
\end{itemize}

\separador
\section{Hoja 5}

\end{document}